% stirling-sec4on5thorder.tex
% Definitions and (corrected) proofs of C. Stirling,
% "An introduction to decidability of higher-order matching", Sections 4--6,
% specialised to interpolation equations of order at most 5.
% HOM project working document, 2026-11-10.
\documentclass[11pt]{article}
\usepackage[margin=2.5cm]{geometry}
\usepackage[T1]{fontenc}
\usepackage{lmodern}
\usepackage{microtype}
\usepackage{amsmath,amssymb,amsthm}
\usepackage{enumerate}
\usepackage{xcolor}
\usepackage[colorlinks=true,linkcolor=blue!50!black,citecolor=blue!50!black]{hyperref}

\newcommand{\Ground}{\mathbf{0}}
\newcommand{\ord}{\mathrm{ord}}
\newcommand{\ar}{\mathrm{ar}}
\newcommand{\Sub}{\mathrm{Sub}}
\newcommand{\depth}{\mathrm{depth}}
\newcommand{\br}{\mathrm{br}}
\newcommand{\Gm}{\mathsf{G}}
\newcommand{\bnd}{\mathrel{\downarrow}}
\newcommand{\suc}[2]{#1{\cdot}#2}
\newcommand{\iv}[2]{\pi[#1,#2]}
\newcommand{\comb}{\mathrm{comb}}
\newcommand{\pth}{\mathrm{path}}
\newcommand{\CE}{C_E}
\newcommand{\sq}{\sqsubseteq}
\newcommand{\ssq}{\sqsubset}
\newcommand{\ra}[1]{\xrightarrow{#1}}
\newcommand{\args}{\mathrm{args}}
\newcommand{\env}{\mathrm{env}}

\newcommand{\Paper}{\textsf{[paper]}}
\newcommand{\Fixed}{\textsf{[corrected]}}
\newcommand{\New}{\textsf{[new]}}
\newcommand{\Sketch}{\textsf{[sketch]}}
\newcommand{\Open}{\textsf{\color{red!60!black}[open]}}
\newcommand{\Coq}{\textsf{[Coq]}}

\theoremstyle{plain}
\newtheorem{theorem}{Theorem}[section]
\newtheorem{proposition}[theorem]{Proposition}
\newtheorem{lemma}[theorem]{Lemma}
\newtheorem{corollary}[theorem]{Corollary}
\newtheorem{fact}[theorem]{Fact}
\newtheorem{claim}[theorem]{Claim}
\theoremstyle{definition}
\newtheorem{definition}[theorem]{Definition}
\newtheorem{example}[theorem]{Example}
\theoremstyle{remark}
\newtheorem{remark}[theorem]{Remark}
\newtheorem*{gapenv}{\gapname}
\newcommand{\gapname}{Gap}
\newcommand{\GapPartition}{\hyperref[gap:ptree]{[Gap-Partition]}}
\newcommand{\GapUnfolding}{\hyperref[gap:unfold]{[Gap-Unfolding]}}
\newcommand{\GapCounting}{\hyperref[gap:count]{[Gap-Counting]}}

\title{Stirling's decidability proof for higher-order matching\\
at order $5$: definitions and corrected proofs}
\author{HOM project working document}
\date{10 November 2026}

\begin{document}
\maketitle

\begin{abstract}
This document restates Sections~4 and~6 of C.~Stirling, \emph{An
introduction to decidability of higher-order matching} (``the paper''),
for interpolation equations of order at most~$5$ only.  Every notion is
specialised to this case: there are two levels, four kinds of bound
variables, four kinds of jumps, and three kinds of stages.  Every result
is either proved, proved in corrected form, or marked as open.  It is
derived from the general working document \texttt{stirling-sec4on.tex}.
Section~5 of the paper (tree automata) is omitted, since the
decidability proof does not use it.
\end{abstract}

\tableofcontents

%======================================================================
\section*{How to read this document}
%======================================================================
\addcontentsline{toc}{section}{How to read this document}

Numbers in brackets, e.g.\ ``[S, Def.~43]'', refer to the paper.  Our
numbering is independent.  Each statement carries a status tag.
\begin{description}
\item[\Paper] Stated and proved as in the paper.  The proof is
  reproduced, possibly with more detail but without change of substance.
\item[\Fixed] The paper's statement or proof is wrong or incomplete as
  written.  We give the corrected statement or proof, and a remark says
  what was changed.
\item[\New] Not in the paper.  Added here, with a proof.
\item[\Sketch] Only a proof sketch is available (in the paper or here).
  The sketch is believed correct, but it has not been checked in full.
\item[\Open] No proof is known.  The statement is used by the paper and
  is listed among the gaps in Section~\ref{sec:summary}.
\item[\Coq] Additionally machine-checked in the Coq development
  \texttt{hom/coq} (Coq~8.18).  In Coq, the results of Sections~4
  and~6.1 are taken as axioms (G1--G10, U1--U7), and everything else is
  proved from them.
\end{description}

\paragraph{Standing restrictions.} As in [S, \S\S2--6.4]:
\begin{enumerate}[(R1)]
\item there is a single base type $\Ground$;
\item the goal $u$ contains no bound variables;
\end{enumerate}
and, throughout this document,
\begin{enumerate}[(R1)]\setcounter{enumi}{2}
\item the interpolation equation has order at most $5$.
\end{enumerate}
Section~\ref{sec:restrictions} summarises how the paper removes (R1) and (R2).

%======================================================================
\section{Background from Sections 2--3}\label{sec:background}
%======================================================================

This section only fixes notation; see the paper for details.

\begin{definition}[Types, order, arity; S, Defs.~1--2]
Types are generated by $A::=\Ground\mid A\to A$.  Every non-base type
is written $(A_1,\dots,A_n,\Ground)$ for $A_1\to\cdots\to A_n\to\Ground$.
\[
\ord(\Ground)=1,\quad \ord(A_1,\dots,A_n,\Ground)=1+\max_i\ord(A_i),
\]
\[
\ar(\Ground)=0,\quad \ar(A_1,\dots,A_n,\Ground)=\max(n,\ar(A_1),\dots,\ar(A_n)).
\]
Note that the base type has order~$1$.  For example, at order $5$ the
level-$2$ variables of a solution have order $\le2$, i.e.\ type
$\Ground\to\cdots\to\Ground\to\Ground$.
\end{definition}

\begin{definition}[Long normal form, interpolation; S, Defs.~6--9]
A term of type $(A_1,\dots,A_n,\Ground)$ is in \emph{lnf} if it is
$\lambda y_1\dots y_n.\,h\,s_1\dots s_m$ with $h$ a variable or a
constant and every $s_j$ in lnf.  An \emph{interpolation equation} $E$
is $X\,w_1\dots w_k=u$ with $X:A=(A_1,\dots,A_k,\Ground)$, closed lnf
terms $w_i:A_i$, and a closed lnf goal $u:\Ground$.  A \emph{solution} is
a closed lnf $t:A$ with $t\,w_1\dots w_k=_{\beta\eta}u$.  The order of
$E$ is $\ord(A)$ and its arity is $\ar(A)$.
\end{definition}

Under (R2) every constant of $u$ is first order, so $u$ is a first-order
ground term.  $\CE$ is the set of constants of $u$ together with a fresh
$b:\Ground$ [S, Def.~10].  $\Sub(u)$ is the set of subterms of $u$, with
subterm order $\sq$.  $\depth(a)=0$ and
$\depth(f\,u_1\dots u_m)=1+\max_i\depth(u_i)$.  $|u|$ is the number of
nodes and $\br(u)$ the number of leaves (branches) of $u$.

\begin{lemma}[S, Prop.~4, Fact~6]\label{lem:reduces}
\Paper{} (1) If $t$ solves $E$ then $t\,w_1\dots w_k\to_\beta^*u$.
(2) If $E$ has a solution, it has one with all constants in $\CE$.
\end{lemma}

\begin{definition}[Trees; S, \S3.1]
The tree of an lnf term uses dummy lambdas.  $\lambda\vec z.\,s$ with
$s:\Ground$ ($\vec z$ possibly empty) is a node labelled $\lambda\vec z$
with one successor.  $h\,s_1\dots s_m$ is a node labelled $h$ with
successors $s_1,\dots,s_m$.  $\suc ni$ is the $i$-th successor of $n$.
The \emph{binding relation} $n\bnd m$ holds when $n$ is labelled
$\lambda z_1\dots z_p$ and binds the variable occurrence at $m$.  Nodes
are \emph{lambda}, \emph{variable} or \emph{constant} nodes; a lambda
node with non-empty $\vec z$ is a \emph{binding node}.

The \emph{interpolation tree} $@\,t\,w_1\dots w_k$ has a root
$\varepsilon$ labelled $@$ with successors $1,\dots,k+1$ carrying the
trees of $t,w_1,\dots,w_k$.
\end{definition}

\subsection{The shape of order-5 equations}\label{sec:shape}

From now on $\ord(A)\le5$.  We fix names for the four kinds of bound
variables.  Recall that $\ord(\Ground)=1$, so a first-order function
$\Ground\to\cdots\to\Ground$ has order~$2$.

\begin{definition}[Types and variable kinds]\label{def:kinds}
Write $A=(A_1,\dots,A_k,\Ground)$ with $\ord(A_i)\le4$,
$A_i=(B_{i1},\dots,B_{im_i},\Ground)$ with $\ord(B_{ih})\le3$, and
$B_{ih}=(C_1,\dots,C_r,\Ground)$ with $\ord(C_g)\le2$, so
$C_g=(\Ground,\dots,\Ground,\Ground)$.
\begin{itemize}
\item In $t=\lambda x_1\dots x_k.\,S$: the \emph{$x$-variables}
  $x_i:A_i$ (order $\le4$), bound by the root of $t$; and the
  \emph{$z$-variables} $z:C_g$ (order $\le2$), bound by the successors
  of $x$-nodes.
\item In $w_i=\lambda y_1\dots y_{m_i}.\,H$: the \emph{$y$-variables}
  $y_h:B_{ih}$ (order $\le3$), bound by the root of $w_i$; and the
  \emph{$v$-variables} $v:\Ground$, bound by the successors of
  $y$-nodes.
\end{itemize}
$q$ bounds $k$, every $m_i$, every $r$, the arity of every $C_g$, and
the arities of the constants of $u$.
\end{definition}

\begin{lemma}[Shape of terms]\label{lem:shape}
\New{} Every potential solution $t$ and every $w_i$ has the form below,
where $\lambda\vec z$, $\lambda\vec v$ may be empty (dummy):
\begin{align*}
S&::=x_i\,(\lambda\vec z.\,S)\dots(\lambda\vec z.\,S)\mid
     z\,(\lambda.\,S)\dots(\lambda.\,S)\mid f\,(\lambda.\,S)\dots(\lambda.\,S)\mid a,\\
H&::=y_h\,(\lambda\vec v.\,H)\dots(\lambda\vec v.\,H)\mid v\mid
     f\,(\lambda.\,H)\dots(\lambda.\,H)\mid a .
\end{align*}
Here $z$ and $v$ range over variables bound above, and the $f$, $a$ in
$t$ are constants of $\CE$.
\end{lemma}
\begin{proof}
The arguments of $x_i$ have the types $B_{ih}$, so they are
$\lambda\vec z.\,S$ with $\vec z:C_1,\dots,C_r$.  The arguments of
$z:C_g$ have base type, so they are dummy lambdas.  The arguments of
$y_h:B_{ih}$ have the types $C_g$, so they are $\lambda\vec v.\,H$ with
$\vec v$ of base type, and $v$-nodes have no successors.  Constants are
first order, so their successors are dummy lambdas.
\end{proof}

\begin{corollary}[Binding nodes and levels]\label{cor:binders}
\New{}
\begin{enumerate}
\item The binding nodes of the interpolation tree are the root of $t$
  (binding $\vec x$), the successors of $x$-nodes (binding $\vec z$),
  the roots of the $w_i$ (binding $\vec y$), and the successors of
  $y$-nodes (binding $\vec v$).  Every other lambda node is a dummy.
\item $x$- and $y$-nodes have level~$1$, $z$- and $v$-nodes have
  level~$2$ (Definition~\ref{def:struct}).  There are no other levels.
\item $z$- and $v$-nodes are end.  An $x$-node is end iff none of the
  $z$-variables bound by its successors occurs; likewise for
  $y$-nodes.
\end{enumerate}
\end{corollary}
\begin{proof}
(1) and (2) are read off Lemma~\ref{lem:shape}.  (3) The successors of
$z$- and $v$-nodes bind nothing.
\end{proof}

\begin{remark}
In the general paper, levels go up to $p$ for order $2p$ or $2p+1$.
Here $p=2$.  If $\ord(A)\le3$, all $\vec z$ are empty and there are no
$z$-variables.
\end{remark}

%======================================================================
\section{Games and interpolation [S, \S4]}\label{sec:games}
%======================================================================

Throughout, $E$ is $X\,w_1\dots w_k=u$ and $t:A$ is a closed lnf term
(a \emph{potential solution}).

\subsection{The game}

\begin{definition}[States, look-up tables; S, Defs.~16--17]\label{def:tables}
$N$ is the set of nodes of $@\,t\,w_1\dots w_k$ and $N_1\subseteq N$ the
set of binding nodes.  The set of \emph{states} is
$S=\{[r]\mid r\in\Sub(u)\}\cup\{[\forall],[\exists]\}$; $[\forall]$ and
$[\exists]$ are \emph{final}.  Look-up tables are defined in layers:
$\Theta_1=\{\varnothing\}$, and $\Theta_{i+1}$ is the set of partial maps
\[
N_1\rightharpoonup \textstyle\bigcup_p N^p\times\bigcup_{j\le i}\Theta_j .
\]
An entry $\theta(m)=((m_1,\dots,m_p),\theta')$ says that the variables
bound by $m$ stand for the subtrees at $m_1,\dots,m_p$, read under the
table $\theta'$.  We write $\args_\theta(m)=(m_1,\dots,m_p)$ and
$\env_\theta(m)=\theta'$.  Two tables are \emph{equal} if they have the
same domain and equal entries; this is well defined because the inner
tables come from earlier positions (Proposition~\ref{prop:tables}(4)).
\end{definition}

\begin{definition}[Moves; S, Def.~19 and Fig.~5]\label{def:moves}
Let the current position be $n[r]\theta$ with $r$ not final.  The next
position depends on the label of $n$.
\begin{description}
\item[$@$:] $\suc n1[r]\theta'$ with
  $\theta'=\theta\{((\suc n2,\dots,\suc n{(k+1)}),\theta)/\suc n1\}$.
\item[$\lambda\vec z$:] $\suc n1[r]\theta$.
\item[$a:\Ground$:] $n[\exists]\theta$ if $r=a$, and $n[\forall]\theta$
  otherwise.
\item[$f$ of arity $p>0$:] if $r=f\,r_1\dots r_p$, then
  $\forall$ chooses $j\in\{1,\dots,p\}$ and the next position is
  $\suc nj[r_j]\theta$; otherwise it is $n[\forall]\theta$.
\item[$z_j:\Ground$:] if $m\bnd n$ and $\theta(m)=((m_1,\dots,m_l),\theta')$
  then $m_j[r]\theta'$.
\item[$z_j$ of type $(B_1,\dots,B_p,\Ground)$, $p>0$:]\ \\ if $m\bnd n$ and
  $\theta(m)=((m_1,\dots,m_l),\theta')$ then $m_j[r]\theta''$ with
  $\theta''=\theta'\{((\suc n1,\dots,\suc np),\theta)/m_j\}$.
\end{description}
The last two moves are \emph{jumps}.
\end{definition}

\begin{definition}[Plays, winning; S, Defs.~18, 20]
A \emph{play} of $\Gm(t,E)$ is a finite sequence
$\pi=\pi_1\dots\pi_l$ of positions $\pi_i=n_i[r_i]\theta_i$ with
$\pi_1=\varepsilon[u]\varnothing$, each $\pi_{i+1}$ a next position of
$\pi_i$, $r_i$ non-final for $i<l$, and $r_l$ final.  $\forall$
\emph{wins} $\pi$ if $r_l=[\forall]$.  She \emph{wins the game} if she
wins some play, and \emph{loses} it otherwise.  We write
$\pi\in\Gm(t,E)$, $\iv ij$ for $\pi_i\dots\pi_j$, and say that $\pi_i$ is
\emph{at} $n_i$.
\end{definition}

\subsection{Basic properties of plays}

\begin{proposition}[S, Prop.~12]\label{prop:tables}
\Paper{} Let $\pi\in\Gm(t,E)$, $j\le l$, and
$\theta_j(m)=((m_1,\dots,m_p),\theta')$.
\begin{enumerate}
\item $\theta_j(n)$ is defined iff $n$ is a binding node above $n_j$.
\item $\theta'(n)$ is defined iff $n$ is a binding node strictly above
  every $m_k$.
\item $m$ is labelled $\lambda z_1\dots z_p$, and $z_k$ has the type of
  $m_k$.
\item $\theta'=\theta_i$ for some $i<j$, and either $i=1$ and
  $m_k=\suc{\varepsilon}{(k+1)}$, or $n_i$ is a variable node and
  $m_k=\suc{n_i}k$.
\end{enumerate}
\end{proposition}
\begin{proof}
Simultaneous induction on $j$.  For $j=1$ the table is empty.  For the
step, look at the move from $\pi_{j-1}$.
\begin{itemize}
\item $@$: $j=2$; the only entry is for node $1$, which is above itself;
  its inner table is $\theta_1=\varnothing$ and nothing lies strictly above
  the roots of the $w_i$.
\item $\lambda$, $a$, $f$: the table is unchanged and $n_j$ is not a
  binding node, or is the successor of a constant, which binds nothing
  since constants of $u$ are first order.
\item jump at $z_q$ bound by $m$, with
  $\theta_{j-1}(m)=((m_1,\dots,m_p),\theta')$: then $n_j=m_q$, and
  $\theta_j=\theta'$ or
  $\theta_j=\theta'\{((\suc{n_{j-1}}1,\dots),\theta_{j-1})/m_q\}$.  By the
  induction hypothesis~(2) for $\theta'$, the binding nodes above $m_q$
  are those strictly above $m_q$ (domain of $\theta'$), plus $m_q$
  itself when $z_q$ has higher type.  This gives (1).  Entries of
  $\theta'$ satisfy (2)--(4) by induction.  The new entry has inner
  table $\theta_{j-1}$, an earlier table, and arguments the successors
  of the variable node $n_{j-1}$.  Its types match because $z_q$ and
  $m_q$ have the same type, by induction~(3).  \qedhere
\end{itemize}
\end{proof}

In particular play never gets stuck: at a variable node the binder has
an entry, and a jump from a variable of type $A$ lands at a lambda node
of type $A$.

\begin{corollary}[The four jumps]\label{cor:jumps}
\New{} At order $\le5$ every jump has one of the following forms.  Let
play be at a variable node $n$ with table $\theta$.
\begin{description}
\item[$x$-jump ($t\to w$).] $n$ carries $x_i$.  Play goes to the root of
  $w_i$, with the single entry $((\suc n1,\dots,\suc n{m_i}),\theta)$
  for it (no entry if $m_i=0$).
\item[$y$-jump ($w\to t$).] $n$ carries $y_h$, bound by the root $\rho$
  of $w_i$, and $\theta(\rho)=((\suc{n'}1,\dots),\theta')$ for an
  $x_i$-node $n'$.  Play goes to $\suc{n'}h$ (a $\lambda\vec z$), with
  table $\theta'$ extended by $((\suc n1,\dots),\theta)$ for
  $\suc{n'}h$.
\item[$z$-jump ($t\to w$).] $n$ carries $z_g$, bound by $\suc{n'}h$, and
  $\theta(\suc{n'}h)=((\suc{n''}1,\dots),\theta')$ for a $y$-node
  $n''$.  Play goes to $\suc{n''}g$ (a $\lambda\vec v$), with table
  $\theta'$ extended by $((\suc n1,\dots),\theta)$ for $\suc{n''}g$
  if $z_g$ has arguments.
\item[$v$-jump ($w\to t$).] $n$ carries $v_e$, bound by $\suc{n''}g$, and
  $\theta(\suc{n''}g)=((\suc{n'''}1,\dots),\theta')$ for a $z$-node
  $n'''$.  Play goes to the dummy lambda $\suc{n'''}e$ with table
  $\theta'$.
\end{description}
So play in $t$ visits the root, then alternates between $t$ and $w$:
$x$- and $z$-jumps leave $t$, $y$- and $v$-jumps return to it.
\end{corollary}
\begin{proof}
Definition~\ref{def:moves}, with the binders of
Corollary~\ref{cor:binders}(1) and Proposition~\ref{prop:tables}(4),
which says which node created each entry.
\end{proof}

\begin{definition}[Associated term; S, proof of Thm.~13]\label{def:assoc}
For a node $n$ and a table $\theta$, $\theta[n]$ is the term
\begin{itemize}
\item $\theta[\suc n1]\,\theta[\suc n2]\dots\theta[\suc n{(k+1)}]$ if $n$ is
  labelled $@$;
\item $a$ if $n$ is labelled $a$, and $f\,\theta[\suc n1]\dots\theta[\suc np]$
  if it is labelled $f$;
\item $\theta[\suc n1]$ for a dummy lambda, and
  $\lambda\vec x.\,\theta[\suc n1]$ for $\lambda\vec x$ with $\vec x$
  non-empty;
\item for $z_j$ bound by $m$: $\theta'[m_j]\,\theta[\suc n1]\dots\theta[\suc nq]$
  if $\theta(m)=((m_1,\dots),\theta')$, and $z_j\,\theta[\suc n1]\dots$
  if $\theta(m)$ is undefined.
\end{itemize}
\end{definition}

\begin{theorem}[S, Thm.~13]\label{thm:game}
\Paper{} $\forall$ loses $\Gm(t,E)$ iff $t$ solves $E$.
\end{theorem}
\begin{proof}
Use \emph{leftmost reduction sequences} with the rules
$(\lambda z_1\dots z_p.s)\,s_1\dots s_p\to s\{s_1/z_1,\dots\}$ and
$f\,s_1\dots s_p\to s_j$ (reduction under a constant, which commits to a
branch).  A term is \emph{terminal} if neither rule applies.  A sequence
$s_1\to\dots\to s_m$ is \emph{compatible} with a binder-free normal
$v:\Ground$ if:
\begin{itemize}
\item for $m=1$: if $s_1$ is terminal then $s_1=v$;
\item for a $\beta$-step: the rest is compatible with $v$;
\item for a step $f\,\vec s\to s_j$: $v=f\,\vec v$ and the rest is
  compatible with $v_j$.
\end{itemize}
A \emph{padding} inserts steps $s\to s$.

\emph{Claim.} For every reduction sequence $\sigma$ from
$t\,w_1\dots w_k$ compatible with $u$ there are a play prefix $\pi$ and
a padding $s_1\to\dots\to s_m$ of $\sigma$ with
$s_i=\theta_{2i-1}[n_{2i-1}]$ for all $i$.  Conversely, every play prefix
has such a sequence.

The proof is by induction on the length, by cases on the label of
$n_{2s-1}$, which is never a lambda node (Fact~\ref{fact:parity}).
\begin{itemize}
\item $@$: $\theta_3[n_3]$ is the body of $t$ with $w_i$ for $y_i$.
\item $f$: $\forall$ follows the branch of the reduction step.
\item $z_j$ of higher type: $\theta_{2s-1}[n_{2s-1}]=\theta'[m_j]\,\vec s$
  is a redex $(\lambda\vec z.s')\vec s$.  The next odd position
  represents $s'\{\vec s/\vec z\}$, because the new entry for $m_j$ maps
  $\vec z$ to the successors of $n_{2s-1}$ under $\theta_{2s-1}$.
\item $z_j:\Ground$: $\theta_{2s-1}[n_{2s-1}]=\theta'[m_j]$ and the play
  moves on without a reduction step; insert a padding step.  Only
  finitely many paddings happen in a row, since each uses a table of
  an earlier position.
\end{itemize}
By strong normalisation, plays are finite.

If $t$ solves $E$, every maximal reduction sequence is finite and
compatible with $u$.  So every play ends with $[\exists]$, because a
play ending with $[\forall]$ would give an incompatible sequence.
Conversely, if $\forall$ loses, then every branch of $u$ is followed by
some play, so by some compatible sequence, and $t\,\vec w$ normalises to
$u$.
\end{proof}

\begin{fact}[S, Fact~10]\label{fact:branches}
\Paper{} If $t$ solves $E$, then $\Gm(t,E)$ has exactly $\br(u)$ plays.
\end{fact}
\begin{proof}
The game is deterministic except for $\forall$'s choices at constants.
If $t$ solves $E$, the state at an $f$-node is always $f\,\vec r$
(otherwise $\forall$ wins), so every choice is available.  Every play
ends at some $a$ with state $[a]$.  Hence plays correspond bijectively
to the branches of $u$.
\end{proof}

\begin{fact}[S, Fact~11]\label{fact:parity}
\Paper{} Let $\pi\in\Gm(t,E)$ have length $l$.
(1) If $i<l$ is even, $n_i$ is a lambda node.
(2) If $i>1$ is odd, $n_i$ is a variable or constant node.
(3) $n_l$ is a constant node and $r_l$ is final.
(4) If $i<j<l$ then $r_j\sq r_i$.
\end{fact}
\begin{proof}
In an lnf tree, lambda nodes and non-lambda nodes alternate along paths.
The $@$-move, the $f$-move and the jumps lead to lambda nodes.  The
$\lambda$-move leads to a non-lambda node.  The $a$-move ends the play
at the same node.  States change only at constants, to a subterm.
\end{proof}

\begin{proposition}[S, Prop.~14]\label{prop:distinct}
\Paper{} If $i\ne j$ and $n_i=n_j$ then $\theta_i\ne\theta_j$.
\end{proposition}
\begin{proof}
If $n_i=n_j$ and $\theta_i=\theta_j$, the associated terms coincide.
By the claim in the proof of Theorem~\ref{thm:game} there is a reduction
sequence from that term back to itself that is not all padding,
contradicting strong normalisation.
\end{proof}

\begin{lemma}[Predecessors]\label{lem:pred}
\New{} If $1<j<l$ and $n_j$ is not a lambda node, then $n_{j-1}$
is the lambda node directly above $n_j$, and $\theta_{j-1}=\theta_j$,
$r_{j-1}=r_j$.
\end{lemma}
\begin{proof}
Only the $\lambda$-move leads to a non-lambda node without ending the
play.
\end{proof}

\begin{lemma}[States change only at constants]\label{lem:states}
\New{} \Coq{} Let $t$ solve $E$ and $i<l$.  Then $r_{i+1}\ne r_i$ iff
$n_i$ is a constant node of positive arity, and then $r_{i+1}\ssq r_i$.
Hence for $i\le j<l$: $r_j\ssq r_i$ iff some $n_h$, $i\le h<j$, is a
constant node of positive arity.
\end{lemma}
\begin{proof}
Only the $f$- and $a$-moves change states; the $a$-move ends the play.
At an $f$-node, $\forall$ loses, so $r_i=f\,\vec r$ and $r_{i+1}=r_j\ssq
r_i$.  The second claim then follows from Fact~\ref{fact:parity}(4).
\end{proof}

\begin{lemma}[Locality]\label{lem:locality}
\New{} \Coq{} Let $t,t'$ be lnf terms of type $A$ and $X$ a
prefix-closed set of nodes on which their interpolation trees agree
(labels, successors, binders).  If no play of $\Gm(t,E)$ has a position
outside $X$, then $\Gm(t',E)=\Gm(t,E)$.
\end{lemma}
\begin{proof}
A move depends only on the current state and table, and on the label,
successors and binder of the current node.  By induction, every
sequence of moves that stays in $X$ is the same in both games.  A
sequence for $t'$ that left $X$ for the first time would, up to that
point, be a sequence for $t$ leaving $X$.
\end{proof}

\subsection{Parent and child positions [S, \S4.3]}

\begin{definition}[Parent; S, Defs.~21--22]\label{def:parent}
Let $\pi\in\Gm(t,E)$.
\begin{enumerate}
\item If $n_j$ is a variable node, $\pi_i$ is the \emph{parent} of
  $\pi_j$ if $n_i\bnd n_j$ and $\theta_i(n_i)=\theta_j(n_i)$.
\item If $n_j$ is a lambda node and $n$ is the node directly above
  $n_j$: if $n=\varepsilon$, the parent is $\pi_1$; if $n$ is a constant
  node, the parent is $\pi_{j-1}$; if $n$ is a variable node, the parent
  is the $\pi_i$ such that $\pi_{i+1}$ is the parent of $\pi_{j-1}$ in the
  sense of item~1.
\end{enumerate}
$\pi_j$ is a \emph{child} of $\pi_i$ if $\pi_i$ is its parent.  In
item~2, $n_j$ is a successor of $n_i$.
\end{definition}

\begin{remark}[Children at order $5$]
By Corollary~\ref{cor:jumps}, a position at an $x$-node has its children
at successors $\lambda\vec z$ of that node, entered by $y$-jumps.  A
position at a $z$-node has its children at the dummy lambdas below it,
entered by $v$-jumps.  Symmetrically in $w$: children of a $y$-position
are entered by $z$-jumps, and $v$-nodes have no successors.  The
children of a constant position are its successors, entered by the
$f$-move.
\end{remark}

\begin{proposition}[S, Prop.~15, Fact~16]\label{prop:parent}
\Paper{} Every position $\pi_j$, $j>1$, at a variable node or at a lambda
node has a unique parent $\pi_i$, and $i<j$.
\end{proposition}
\begin{proof}
Variable node $n_j$ with binder $m$: an entry for $m$ enters a table
only when play is at $m$ (the $@$-move or a higher-type jump), so some
earlier $\pi_i$ at $m$ has $\theta_i(m)=\theta_j(m)$.  If there were two,
$\pi_i$ and $\pi_{i'}$, then $\pi_{i-1}$ and $\pi_{i'-1}$ would be at the
same node with the same table (the entry records both), contradicting
Proposition~\ref{prop:distinct}.  Lambda node: use the three cases of
Definition~\ref{def:parent}(2) and the variable case.
\end{proof}

\begin{definition}[Chain; S, Def.~23]\label{def:chain}
For $\pi_j$ at a lambda node, a \emph{chain} for $\pi_j$ is a sequence
$\pi_{j_1},\dots,\pi_{j_{2p}}$ with
(1) $\pi_{j_{2p}}=\pi_j$;
(2) $\pi_{j_{2m-1}}$ is the parent of $\pi_{j_{2m}}$ for $1\le m\le p$;
(3) $j_{2m+1}=j_{2m}+1$ for $1\le m<p$.
The nodes $n_{j_1},\dots,n_{j_{2p}}$ then form a downward path.
\end{definition}

\begin{definition}[S, Def.~24]
$\theta\le\theta'$ ($\theta'$ \emph{extends} $\theta$) if $\theta'(m)=\theta(m)$
whenever $\theta(m)$ is defined.
\end{definition}

\begin{proposition}[S, Props.~17--18]\label{prop:mono}
\Paper{}
(1) If $n_j$ is a lambda node with parent $\pi_i$, then
$\theta_i\le\theta_j$.
(2) For a chain $\pi_{j_1},\dots,\pi_{j_{2p}}$:
$\theta_{j_1}\le\cdots\le\theta_{j_{2p}}$, and if $n_{j_a}\bnd n_{j_b}$
then $\pi_{j_a}$ is the parent of $\pi_{j_b}$.
\end{proposition}
\begin{proof}
(1) Cases on the label of $n_i$.
\begin{itemize}
\item $@$: $\theta_1=\varnothing$.
\item Constant: $\theta_j=\theta_i$.
\item Variable $x$: $\pi_{i+1}$ is the parent of $\pi_{j-1}$ (at $z_k$),
  so $\theta_{j-1}(n_{i+1})=\theta_{i+1}(n_{i+1})=((\dots),\theta_i)$.
  Hence $\theta_j=\theta_i$ or $\theta_i\{\dots/n_j\}$, and $\theta_i$
  has no entry for $n_j$ (Proposition~\ref{prop:tables}(1)).
\end{itemize}
(2) Chain links are parent/child pairs, and the step $j_{2m}\to
j_{2m}+1$ is a $\lambda$-move, which keeps the table.  The second claim
follows with Proposition~\ref{prop:parent}.
\end{proof}

\begin{proposition}[S, Prop.~19]\label{prop:agree}
\Paper{} If $\theta_i(m)=\theta_j(m)$ are defined and $m'$ is a binding
node above $m$, then $\theta_i(m')=\theta_j(m')$.
\end{proposition}
\begin{proof}
Take the chains for $\pi_i$ and $\pi_j$ from the node above $m'$, and
their prefixes ending at $m$.  The last positions of the prefixes have
equal entries for $m$.  So they coincide, since otherwise their
predecessors would be at the same node with the same table
(Proposition~\ref{prop:distinct}).  Hence the prefixes coincide, and
monotonicity gives the claim.
\end{proof}

\begin{lemma}[Ancestors are visited first]\label{lem:ancestors}
\New{} \Coq{} If $\pi_j$ is at $n$ and $m\ne\varepsilon$ is a
non-lambda node strictly above $n$, then some $i<j$ has $n_i=m$.  If $n$
is a lambda node, there is a unique chain for $\pi_j$ whose first
position is at $m$.
\end{lemma}
\begin{proof}
Let $\rho(j)=j-1$ at non-lambda nodes (Lemma~\ref{lem:pred}) and
$\rho(j)=$ the parent index at lambda nodes.  Then $n_{\rho(j)}$ is the
node directly above $n_j$, and iterating $\rho$ walks up the path.
The visited positions form a chain, and $\rho$ is a function, which
gives uniqueness.
\end{proof}

\begin{lemma}[Alternation]\label{lem:alt}
\New{} If $n_j$ is a variable node then $n_{j+1}$ is a lambda node, and
exactly one of $n_j,n_{j+1}$ lies in the $t$-component.
\end{lemma}
\begin{proof}
By Proposition~\ref{prop:tables}(4), every entry for a binder in $w_q$
points into $t$ and vice versa, since every $w_q$ is closed and $t$ is
closed.
\end{proof}

%======================================================================
\section{Uniformity of game playing [S, \S6.1]}\label{sec:uniformity}
%======================================================================

\begin{proposition}[S, Prop.~23]\label{prop:between}
\Paper{} If $j<m$ and $\pi_i$ is the parent of both $\pi_j$ and $\pi_m$,
then $\pi_j$ is the parent of some $\pi_k$ with $j<k<m$.
\end{proposition}
\begin{proof}
\emph{Variable case.} $n_i\bnd n_j,n_m$ and
$\theta_j(n_i)=\theta_m(n_i)=((m_1,\dots,m_p),\theta_l)$.  Then $\pi_{j+1}$ is
at some $m_q$ with $\theta_{j+1}=\theta_l\{((\suc{n_j}1,\dots),\theta_j)/m_q\}$.
$\theta_l$ has no entry for $n_i$, since the jump crosses between $t$ and
some $w$.  An entry of $\theta_l$, or of a table reachable from it,
cannot be the entry of $n_i$ recorded at $\pi_i$, by
Proposition~\ref{prop:distinct}.  So the only way back to a table
containing that entry is through a variable bound by $m_q$, i.e.\
through a child of $\pi_j$.

\emph{Lambda case.} Cases on the node above.  For $@$, play leaves $t$
(resp.\ $w_q$) by a jump from a child of $\pi_j$.  A constant node has
only one child position.  For a variable node, apply the variable case
to $\pi_{i+1}$, $\pi_{j-1}$, $\pi_{m-1}$.
\end{proof}

\begin{definition}[Excursion]\label{def:excursion}
\New{} Let $\pi_j$ be at a variable node $n_j$, $j<l$.  The
\emph{excursion} at $\pi_j$ is the interval $\iv{j+1}{j'-1}$, where
$\pi_{j'}$ is the first position after $\pi_j$ that is a child of
$\pi_j$; if there is none, it is $\iv{j+1}l$.  So the excursion is the
part of the play between the jump out of $n_j$ and the first return to a
successor of $n_j$ with the table of $\pi_j$.  The excursion
\emph{returns} if $\pi_{j'}$ exists, and then it returns \emph{to}
$\suc{n_j}h$ if $n_{j'}=\suc{n_j}h$.  It \emph{contributes} if
$r_{j+1}\ne r_{j'-1}$ (resp.\ $r_l$).
\end{definition}

\begin{remark}
At order~$5$ (Corollary~\ref{cor:jumps}): the excursion at an
$x$-node of $t$ runs from the root of some $w_i$ until a $y$-jump back
to a successor $\lambda\vec z$ of that node; the excursion at a
$z$-node runs from a $\lambda\vec v$ in some $w_i$ until a $v$-jump back
to a dummy lambda below that node.  Proposition~\ref{prop:24} compares
initial segments of two excursions; the intervals in
Definition~\ref{def:compl} are such segments.
\end{remark}

\begin{definition}[Correspondence; S, Def.~32]
$\iv i{i+m}$ and $\iv j{j+m}$ \emph{correspond} if $n_{i+h}=n_{j+h}$ for
$0\le h\le m$.
\end{definition}

\begin{proposition}[S, Prop.~24]\label{prop:24}
\Paper{} Let $\pi_i$ be the parent of $\pi_j$ and of $\pi_m$, $j<m$, and
let $n_j,n_m$ carry the same variable.  Then there is $k$ such that
$\iv{j+1}{j+k}$ and $\iv{m+1}{m+k}$ correspond and either
\begin{enumerate}
\item[(1)] $\pi_j$ is the parent of $\pi_{j+k+1}$, $\pi_m$ is the parent of
  $\pi_{m+k+1}$, and $n_{j+k+1}=\suc{n_j}p$ iff $n_{m+k+1}=\suc{n_m}p$; or
\item[(2)] $n_{j+k}$ is a constant node, and $\pi_{m+k}$ is final or
  $n_{j+k+1}\ne n_{m+k+1}$.
\end{enumerate}
\end{proposition}
\begin{proof}
Write $\theta_i(n_i)=((m_1,\dots),\theta_l)$ and let $n_j,n_m$ carry $z_q$.
Then $n_{j+1}=n_{m+1}=m_q$ and
\[
\theta_{j+1}=\theta_l\{((\suc{n_j}1,\dots),\theta_j)/m_q\},\qquad
\theta_{m+1}=\theta_l\{((\suc{n_m}1,\dots),\theta_m)/m_q\}.
\]
These differ only in the entry for $m_q$.  Hence the two continuations
make the same moves until play reaches either a variable bound by
$m_q$, which gives case~(1) with the same successor index, or a
constant where $\forall$'s choices or the outcomes differ, which gives
case~(2).  This is the special case of Proposition~\ref{prop:24star} in
which the only difference is at $m_q$.  The hypothesis $n_j\ne n_m$ of
the paper is not needed.
\end{proof}

\subsection{A generalisation without common parent}

\begin{definition}[Difference pairs]\label{def:diff}
\New{} For tables $\theta,\theta'$, let $P(\theta,\theta')$ be the least
set of pairs of tables such that
\begin{itemize}
\item $(\theta,\theta')\in P$;
\item if $(\rho,\rho')\in P$ and $\args_\rho(m)=\args_{\rho'}(m)$, then
  $(\env_\rho(m),\env_{\rho'}(m))\in P$.
\end{itemize}
A \emph{difference} of $(\theta,\theta')$ is a triple $(\rho,\rho',m)$ with
$(\rho,\rho')\in P$ and $\args_\rho(m)\ne\args_{\rho'}(m)$.  We write
$\theta\equiv\theta'$ if there is no difference.  Under the reading of
tables as substitutions (Definition~\ref{def:assoc}), $\theta\equiv\theta'$
means that the two tables substitute the same subtrees, hereditarily.
\end{definition}

\begin{proposition}[Prop.~24\textsuperscript{*}]\label{prop:24star}
\New{} Let $a<b$ and $n_a=n_b$.  Then there is $k\ge0$ such that
$\iv a{a+k}$ and $\iv b{b+k}$ correspond and one of the following
holds.
\begin{enumerate}
\item[(1)] \emph{(binder split)} $n_{a+k}$ is a variable $z_q$ bound by
  $m$, and the entries $\theta_{a+k}(m)$, $\theta_{b+k}(m)$ are
  $\rho(m),\rho'(m)$ for a difference $(\rho,\rho',m)$ of
  $(\theta_a,\theta_b)$.  Then $\pi_{a+k+1}$ and $\pi_{b+k+1}$ are at the
  $q$-th successors of the nodes that created these entries.  They are
  children (Definition~\ref{def:parent}) of the creating positions.
\item[(2)] \emph{($\forall$ split)} $n_{a+k}$ is a constant node, and
  $\pi_{a+k}$ or $\pi_{b+k}$ is final, or $\forall$ takes different
  branches.
\end{enumerate}
In particular, if $\theta_a\equiv\theta_b$ then only case~(2) can occur.
\end{proposition}
\begin{proof}
Let $\hat P$ be the closure of $P(\theta_a,\theta_b)$ under the
following operation: from $(\rho,\rho')\in\hat P$ and a pair
$(\sigma,\sigma')\in\hat P$, form
$(\rho\{((\vec n),\sigma)/m\},\rho'\{((\vec n),\sigma')/m\})$, with the
same $\vec n$ and $m$ on both sides.  By induction on $\hat P$: every
entry pair of a table pair in $\hat P$ either has equal args (and an env
pair in $\hat P$), or is an entry pair of some table pair in $P$.

\emph{Invariant:} as long as the runs correspond,
$(\theta_{a+i},\theta_{b+i})\in\hat P$.  Check it move by move
(Definition~\ref{def:moves}):
\begin{itemize}
\item At $\lambda$ and constant nodes the table is unchanged.  At a
  constant node only the state or $\forall$'s choice can separate the
  runs; this is case~(2).
\item At $@$ (only at $a=b=1$, impossible since $a<b$) or at a
  higher-type jump, the new entry has $\vec n$ the successors of the
  common node and env the current pair, so the invariant is kept.
\item At a variable $z_q$ bound by $m$:
  \begin{itemize}
  \item if $\args(m)$ agree, the next nodes agree, and the new tables
    come from the env pair, which lies in $\hat P$;
  \item if they differ, the entry pair is not a created one, so it is
    an entry pair of a pair in $P$.  This is case~(1).
  \end{itemize}
\end{itemize}
The successor and child claims follow from
Proposition~\ref{prop:tables}(4) and Definition~\ref{def:parent}.
Plays are finite, so $k$ exists.
\end{proof}

\begin{corollary}[24$'$: no common parent]\label{cor:24prime}
\New{} Let $j<m$ and let $n_j,n_m$ carry the same variable $z_q$ of a
binder $b$.  Let the parents of $\pi_j,\pi_m$ be visits $\pi_i,\pi_{i'}$
of $b$ with $\theta_i(b)=((\vec m),\theta_l)$ and
$\theta_{i'}(b)=((\vec m),\theta_{l'})$ (same arguments).
\begin{itemize}
\item If $\theta_l\equiv\theta_{l'}$, the conclusion of
  Proposition~\ref{prop:24} holds.
\item Otherwise there is a third outcome: a binder split
  (Proposition~\ref{prop:24star}(1)) at a difference of
  $(\theta_l,\theta_{l'})$.
\end{itemize}
\end{corollary}
\begin{proof}
Apply Proposition~\ref{prop:24star} to $a=j+1$, $b=m+1$.  The
differences of $(\theta_{j+1},\theta_{m+1})$ are the one at $m_q$ (if
$n_j\ne n_m$) together with those of $(\theta_l,\theta_{l'})$.
\end{proof}

\subsection{Complementary intervals and structure of terms}

\begin{definition}[Complementary; S, Def.~33]\label{def:compl}
Let $j<m$ and let $n_j,n_m$ carry the same variable.  The intervals
$\iv{j+1}{j+k}$ and $\iv{m+1}{m+k}$ are \emph{complementary} if case~(1)
or case~(2) of Proposition~\ref{prop:24} holds for them.  No common
parent is required.
\end{definition}

\begin{definition}[Levels, embedded, families; S, Defs.~34--36]\label{def:struct}
\begin{itemize}
\item A variable node has \emph{level}~$1$ if it is bound by the root
  of $t$ or of some $w_q$ ($x$- and $y$-nodes), and level~$2$ if it is
  bound by a successor of a level-$1$ node ($z$- and $v$-nodes).
\item A node $m$ is \emph{embedded} if some $n\ne m$ above $m$ carries
  the same variable and has the same binder.
\item $n$ is a \emph{descendant under binding} of $m$ if $n=m$, or some
  successor $m'$ of $m$ binds some $n'$ and $n$ is a descendant under
  binding of $n'$.
\item Two nodes are in the same \emph{family} if they have a common
  ancestor under binding.
\item A node is \emph{end} if it has no descendant under binding other
  than itself.
\end{itemize}
\end{definition}

\begin{remark}[Families at order $5$]
For an $x$-node $n$ let $F(n)$ consist of $n$ and the $z$-nodes bound by
successors of $n$.  These are exactly the descendants under binding of
$n$, and a $z$-node has no descendant but itself.  Taking ancestors
among variable nodes of $t$, the families of $t$ are therefore the sets
$F(n)$.  In particular every family has a unique $x$-node, its
\emph{head}, and two $x$-nodes are never in the same family.
\end{remark}

\begin{lemma}[Levels]\label{lem:levels}
\New{} Every variable node has level $1$ or $2$: $x$- and $y$-nodes
have level~$1$, and $z$- and $v$-nodes level~$2$.  Level-$1$ variables
of $t$ have order $\le4$, level-$2$ variables of $t$ order $\le2$.
\end{lemma}
\begin{proof}
Corollary~\ref{cor:binders} and Definition~\ref{def:kinds}.
\end{proof}

\begin{fact}[S, Fact~25]\label{fact:25}
\Paper{} If $n_i$ is end, at most one position $\pi_j$ with $j>i$ is a
child of $\pi_i$.
\end{fact}
\begin{proof}
Suppose $\pi_j$ and $\pi_m$ ($j<m$) are both children of $\pi_i$.  They
are at successors of $n_i$, and by Definition~\ref{def:parent}(2)
$\pi_{i+1}$ is the parent of both $\pi_{j-1}$ and $\pi_{m-1}$.  By
Proposition~\ref{prop:between}, $\pi_{j-1}$ is the parent of some
$\pi_k$ with $j-1<k<m-1$.  By Definition~\ref{def:parent}(2) again, this
means that $\pi_j$ is the parent of $\pi_{k-1}$.  So the successor $n_j$
of $n_i$ binds the variable at $n_{k-1}$, and $n_i$ is not end.
\end{proof}

\begin{fact}[S, Fact~26]\label{fact:26}
\Paper{} If $n_j,n_m$ ($j<m$) are $x$-nodes carrying the same
variable, then for some $k$ the intervals $\iv{j+1}{j+k}$ and
$\iv{m+1}{m+k}$ are complementary.
\end{fact}
\begin{proof}
Both are bound by the root of $t$, which is visited only at $\pi_2$.  So
$\pi_2$ is their common parent; apply Proposition~\ref{prop:24}.
\end{proof}

\begin{fact}[S, Fact~27]\label{fact:27}
\Sketch{} Let $n,m$ carry the same variable with the same binder $b$,
let $m$ lie below $\suc np$, and let $n_j=m$.  Then there are $i<j$ and
$k$ such that
\begin{enumerate}
\item $n_i=n$ and $\iv{i+1}{i+k}$, $\iv{j+1}{j+k}$ are complementary;
\item if $n,m$ are end and $r_{i+1}=r_{i+k}$, then $n_{j+k+1}=\suc{n_j}p$.
\end{enumerate}
\end{fact}
\begin{proof}[Proof sketch]
(1) By Lemma~\ref{lem:ancestors} the chain for $\pi_{j-1}$ from the
position at $b$ passes a position $\pi_i$ at $n$.  By
Proposition~\ref{prop:mono} it has the same entry for $b$, hence the
same parent as $\pi_j$.  Apply Proposition~\ref{prop:24}.

(2) The chain passes a child of $\pi_i$ at $\suc np$.  If $n$ is end,
that is the only child (Fact~\ref{fact:25}).  In case~(1) of
Proposition~\ref{prop:24} the return is to that child, so the index is
$p$.  Excluding case~(2) needs the non-contribution hypothesis; the
paper does not spell this out.
\end{proof}

\begin{proposition}[S, Prop.~28]\label{prop:28}
\Paper{} Let $\pi_j$ be a position in $w_q$ and
$\pi_{j_1},\dots,\pi_{j_{2p}}$ a chain for $\pi_j$.  If
$\pi_{j_{2a-1}},\pi_{j_{2b-1}}$ ($a\le b$) are positions in $t$, then
$n_{j_{2a-1}}$ and $n_{j_{2b-1}}$ belong to the same family.
\end{proposition}
\begin{proof}
Both are variable nodes, since their children lie in $w_q$.  Extend the
chain upward to the successor of the root of $w_q$.  The position before
the root of $w_q$ is at an $x$-node of $t$.  That node is an
ancestor under binding of every $t$-node on the chain, since each jump
back into $w_q$ is through a variable bound inside the chain path.
\end{proof}

\begin{definition}[Recurrence property; S, Def.~37]\label{def:recurrence}
A chain $\pi_{j_1},\dots,\pi_{j_{2p}}$ has the \emph{recurrence property}
if, for every $i\le p$ with $n_{j_{2i-1}}$ a variable node, there are
$j<j_{2i-1}$ and $k$ such that $\iv{j+1}{j+k}$ and
$\iv{m'+1}{m'+k}$ are complementary, where $m'=j_{2i-1}$ and
$j_{2i}=m'+k+1$.
\end{definition}

\begin{proposition}[S, Prop.~29]\label{prop:29}
\Fixed{} Let $n_i$ be a variable or constant node in some $w_q$.  There
is a largest $r\ge0$ such that the \emph{greedy} chain
$\pi_{j_1},\dots,\pi_{j_{2r}}$ with $j_1=i$ has the recurrence property.
Moreover $r$ is at most the depth of $w_q$.
\end{proposition}
\begin{proof}
The greedy chain is built as follows.  Given a recurrent greedy chain of
length $2r'$, let $\pi_{i'}$ be $\pi_i$ (if $r'=0$) or
$\pi_{j_{2r'}+1}$, and let $\pi_{j'}$ be the \emph{first} later child of
$\pi_{i'}$, if any.
\begin{itemize}
\item At a constant node, $j'=i'+1$, and the chain extends.
\item At a variable node, it extends iff for some earlier position at a
  node with the same variable the excursions (Definition~\ref{def:excursion}) $\iv{\cdot+1}{\cdot+k}$ and
  $\iv{i'+1}{i'+k}$, $k=j'-i'-1$, are complementary.
\end{itemize}
Otherwise the construction stops.  The nodes form a downward path in
$w_q$, which bounds $r$.
\end{proof}

\begin{remark}
The paper states only ``there is a largest $p$''.  Stages
(Definition~\ref{def:stages}) need one specific chain.  We take the
greedy one, which is what the paper's proof constructs.  A recurrent
chain that uses a later child is not excluded by Definition~\ref{def:recurrence}
alone, so the choice matters.
\end{remark}

%======================================================================
\section{Two transformations [S, \S6.2]}\label{sec:transformations}
%======================================================================

\begin{definition}[Replacement, size; S, \S6.2]
\Fixed{} For a variable or constant node $n$ of $t$ and $s:\Ground$,
$t[s/n]$ replaces the subtree at $n$ by $s$.  The size measure is
$\mu(t)=(|t|,\nu(t),\kappa(t))$, ordered lexicographically:
\begin{itemize}
\item $|t|$ is the number of nodes, dummy lambdas included;
\item $\nu(t)$ is the number of variable nodes;
\item $\kappa(t)$ is the number of nodes labelled by a constant other
  than $b$.
\end{itemize}
The measure is used only to show that $T_1$ and $T_2$ terminate
(Lemma~\ref{lem:terse}).
\end{definition}

\begin{definition}[$T_1$; S, Def.~38]\label{def:T1}
\Fixed{} If $n$ is a variable node or a constant node other than $b$,
and no play of $\Gm(t,E)$ is ever at $n$, then $t\,T_1\,t[b/n]$.
\end{definition}

\begin{remark}
The paper applies $T_1$ only to variables and to constants of positive
arity, and measures size by $(|t|,\nu(t))$.  A terse solution
(Definition~\ref{def:terse}) may then keep unvisited leaves $\lambda.\,a$ with $a\ne b$ of arity~$0$.
Adding the 0-ary constants other than $b$, and the third component
$\kappa$, removes them.  This gives the clean statement of
Corollary~\ref{cor:visited}.  The rest of the proof only uses $\mu$
through ``$T_1$ and $T_2$ make $t$ smaller'' (Lemma~\ref{lem:terse}),
so nothing else changes.
\end{remark}

\begin{proposition}[S, Fact~30]\label{prop:T1}
\Fixed{} \Coq{} If $t$ solves $E$ and $t\,T_1\,t'$, then $t'$ solves $E$
and $\mu(t')<\mu(t)$.
\end{proposition}
\begin{proof}
By Lemma~\ref{lem:ancestors} no play visits a node below $n$ either.
Let $X$ be the set of nodes not below $n$.  Lemma~\ref{lem:locality}
gives $\Gm(t',E)=\Gm(t,E)$, and Theorem~\ref{thm:game} concludes.  For
the size:
\begin{itemize}
\item if the removed subtree has more than one node, $|t|$ drops;
\item if it is a single variable $z:\Ground$, $|t|$ stays and $\nu$
  drops;
\item if it is a single constant $a\ne b$ of arity $0$, $|t|$ and $\nu$
  stay and $\kappa$ drops.
\end{itemize}
The Coq proof (\texttt{T1\_sound}) covers the paper's version of $T_1$;
the extension to 0-ary constants is not yet in Coq.
\end{proof}

\begin{lemma}[$b$ is never visited]\label{lem:bnotvisited}
\New{} If $t$ solves $E$, no play of $\Gm(t,E)$ is at a node labelled
$b$, nor at the dummy lambda directly above such a node.
\end{lemma}
\begin{proof}
$b$ does not occur in $u$.  A position at a $b$-node has a state $[r]$
with $r\in\Sub(u)$, so $r\ne b$, and the $a$-move gives $[\forall]$,
contradicting Theorem~\ref{thm:game}.  A position at the lambda above
is always followed by a position at its unique successor, the $b$-node.
\end{proof}

\begin{corollary}[Visited nodes]\label{cor:visited}
\New{} Let $t$ solve $E$, and suppose $T_1$ does not apply to $t$ (for
instance, $t$ is terse, Definition~\ref{def:terse}).  A node of $t$ is visited by
some play iff it is neither a $b$-node nor the dummy lambda directly
above a $b$-node.
\end{corollary}
\begin{proof}
($\Rightarrow$) Lemma~\ref{lem:bnotvisited}.

($\Leftarrow$) Let $n$ be unvisited.
\begin{itemize}
\item If $n$ is a variable node or a constant node other than $b$, then
  $T_1$ applies to $n$, contrary to the assumption.
\item Otherwise $n$ is a lambda node.  Its successor $n'$ is not visited
  either, since a position at $n'$ is preceded by one at $n$
  (Lemma~\ref{lem:pred}).  By the first case $n'$ is a $b$-node.
\end{itemize}
The root of $t$ is visited at $\pi_2$, so these are all the cases.
\end{proof}

\begin{lemma}[Disjoint contributions]\label{lem:disjoint}
\New{} \Coq{} Let $t$ solve $E$.  If
$\iv{i_1}{j_1},\dots,\iv{i_s}{j_s}$ are intervals of one play with
$r_{i_h}\ne r_{j_h}$ and $j_h\le i_{h+1}$, then $s\le\depth(u)$.
\end{lemma}
\begin{proof}
$r_{i_1}\sqsupset r_{j_1}\sqsupseteq r_{i_2}\sqsupset\cdots$ is a chain in
$\Sub(u)$.
\end{proof}

\begin{lemma}[Constants on a path]\label{lem:constants}
\New{} Let $t$ solve $E$.  On every path of $t$, at most $\depth(u)$
constant nodes of positive arity are visited by some play.
\end{lemma}
\begin{proof}
Take the lowest such constant $c$ and a play that passes from $c$ to its
successor on the path; such a play exists by Fact~\ref{fact:branches}.
The chain from the root (Lemma~\ref{lem:ancestors}) passes every
constant above it with an $f$-move, and each $f$-move strictly
decreases the state.
\end{proof}

\begin{lemma}[Contributing moves over all plays]\label{lem:global}
\New{} Let $t$ solve $E$ and identify positions of plays with a common
prefix.  There are at most $|u|-\br(u)$ distinct $f$-moves in the tree
of plays.
\end{lemma}
\begin{proof}
An $f$-move at state $[r]$ consumes the internal node of $u$ at which
$r$ occurs.  Here states are read as occurrences, determined by
$\forall$'s choices.  Two plays with the same choices up to that point
coincide up to that point.
\end{proof}

\begin{definition}[End path; S, Def.~39]\label{def:endpath}
\Fixed{} A downward path $\gamma=m_1,\dots,m_{2p}$ ($p\ge1$) of $t$ is
\emph{end} if
\begin{enumerate}
\item $m_1$ is a variable node;
\item for $1<i\le p$, $m_{2i-1}$ is a variable node whose binder lies on
  $\gamma$;
\item every variable node strictly below $m_{2p}$ is bound by a node
  strictly above $m_1$ or \emph{strictly} below $m_{2p}$.
\end{enumerate}
\end{definition}

\begin{remark}
The paper says ``below $m_{2p}$''.  It must be read strictly: if
$m_{2p}$ bound a variable below it, that variable would become free after
$T_2$.
\end{remark}

\begin{lemma}[End paths at order $5$]\label{lem:endpath5}
\New{} The end paths of $t$ are exactly the following.
\begin{enumerate}
\item \emph{$z$-paths}: $m_1$ a $z$-node and $m_2$ one of its dummy
  successors ($p=1$).  Every such pair is an end path.
\item \emph{$x$-paths}: $m_1$ an $x$-node $n$, $m_2=\suc nh$, and, if
  $p\ge2$, $m_3,m_5,\dots,m_{2p-1}$ are $z$-nodes bound by $m_2$, each
  directly below the dummy lambda $m_{2i}$ under the previous one.  The
  condition is that no variable strictly below $m_{2p}$ is bound by
  $m_2$.
\end{enumerate}
\end{lemma}
\begin{proof}
$m_1$ is a variable node of $t$, so a $z$- or $x$-node.  The binder of
$m_3$ lies on $\gamma$ above $m_3$, so it is $m_2$.  If $m_1$ is a
$z$-node, $m_2$ is a dummy, so $p=1$; condition~(3) holds, since a
variable node of $t$ strictly below $m_2$ is an $x$-node (bound at the
root) or a $z$-node, whose binder is a successor of an $x$-node and so is
not $m_2$.  If $m_1=n$ is an $x$-node, every $m_{2i-1}$, $i\ge2$, is bound
by a binding node on $\gamma$; the even nodes after $m_2$ are dummies
below $z$-nodes, so that binder is $m_2$, and the $m_{2i-1}$ are
$z$-nodes of $F(n)$.  Condition~(3) then says exactly that nothing
strictly below $m_{2p}$ is bound by $m_2$, since $m_2$ is the only
binder on $\gamma$.
\end{proof}

\begin{definition}[Contribution; S, Def.~40]\label{def:contrib}
\Fixed{} An interval $\iv ij$ with $j<l$ \emph{contributes} if
$r_i\ne r_j$.  The paper requires $n_j$ to be a lambda node; we drop
this, so that a final stage $\iv ij$ contributes iff $\iv i{j-1}$ does,
as [S, Def.~44] needs.
\end{definition}

\begin{definition}[Redundancy; S, Def.~41]\label{def:redundant}
Let $\gamma$ be an end path.  $\gamma$ is \emph{redundant in} $\iv ik$ if
either no $j_1\in[i,k)$ has $n_{j_1}=m_1$, or, for the least such $j_1$,
\begin{enumerate}[(a)]
\item there are $j\le k$ and a chain $\pi_{j_1},\dots,\pi_{j_{2p}}$ for
  $\pi_j$ with $n_{j_h}=m_h$ for all $h$;
\item $\iv{j_1}j$ does not contribute;
\item $\gamma$ is redundant in $\iv jk$.
\end{enumerate}
$\gamma$ is \emph{redundant in} $\Gm(t,E)$ if it is redundant in $\iv1l$
for every play.
\end{definition}

\begin{definition}[$T_2$; S, Def.~42]
If $\gamma$ is an end path redundant in $\Gm(t,E)$ and $t'$ is the
subtree at $\suc{m_{2p}}1$, then $t\,T_2\,t[t'/m_1]$.  By
Definition~\ref{def:endpath}(3) the result is closed, and $|t|$ drops.
\end{definition}

\begin{lemma}[$T_2$ simulation]\label{lem:T2sim}
\New{} Let $t$ solve $E$, let $\gamma=m_1,\dots,m_{2p}$ be an end path
redundant in $\Gm(t,E)$, let $s=\suc{m_{2p}}1$, let $S$ be the subtree of
$t$ at $s$, and let $t''=t[S/m_1]$.  Put
\[
U=\{\text{nodes of the subtree of $t$ at $m_1$}\}\setminus S ,
\]
so $\gamma\subseteq U$, and let $\Gamma$ be the set of binding nodes on
$\gamma$.  Let $\varrho$ be the bijection from the nodes of
$@\,t''\,w_1\dots w_k$ onto $N\setminus U$ with
$\varrho(m_1w)=s\,w$ and $\varrho(n)=n$ for $n$ not below $m_1$.
\begin{enumerate}
\item \emph{(Decomposition.)} Every $\pi\in\Gm(t,E)$ has the form
  \[
  \pi=\sigma_1\,D_1\,\sigma_2\cdots D_q\,\sigma_{q+1},\qquad
  D_r=\iv{a_r}{j_r},
  \]
  where $n_{a_r}=m_1$, $n_{j_r}=m_{2p}$, there is a chain for
  $\pi_{j_r}$ starting at $\pi_{a_r}$ and running along $\gamma$,
  $r_{a_r}=r_{j_r}$, no position of a $\sigma_r$ is at $m_1$, and every
  $\sigma_r$ is non-empty.  No $D_r$ contains an $f$-move.
\item \emph{(Invariant.)} Call a position of $\pi$ \emph{outer} if it
  lies in some $\sigma_r$.  Call a table $\theta$ \emph{clean} if every
  entry $\theta(b)$ with $b\notin\Gamma$ has the form
  $((\suc v1,\dots,\suc vh),\theta_i)$ with $\pi_i$ an outer position.
  Then every outer position $\pi_i$ satisfies $n_i\notin U$, and its
  table $\theta_i$ is clean.
\item \emph{(Simulation.)} For a clean table $\theta$ let $\theta^\flat$
  be defined by recursion on layers: its domain is
  $\varrho^{-1}(\operatorname{dom}\theta\setminus\Gamma)$, and
  $\theta^\flat(\varrho^{-1}b)=((\varrho^{-1}\suc v1,\dots),\theta_i^\flat)$
  when $\theta(b)=((\suc v1,\dots),\theta_i)$.  For an outer position put
  $\varphi(n[r]\theta)=\varrho^{-1}(n)[r]\theta^\flat$.  Then the images
  of the outer positions of $\pi$, in order, form a play $\varphi(\pi)$
  of $\Gm(t'',E)$ with the same final state, and $\forall$ makes the same
  choices in $\pi$ and in $\varphi(\pi)$.
\item \emph{(Converse.)} Every play of $\Gm(t'',E)$ is $\varphi(\pi)$ for
  some $\pi\in\Gm(t,E)$.
\end{enumerate}
\end{lemma}
\begin{proof}
Write $\lambda_0$ for the lambda node directly above $m_1$; it exists,
since $m_1$ is a variable node, and it is not in $U$.

\emph{Preliminaries on $\varrho$.}  $t''$ is $t$ with the subtree at
$m_1$ replaced by $S$.  Both $m_1$ and $s$ are non-lambda nodes, so of
ground type, and $t''$ is in lnf.  $\varrho$ preserves labels and
successors.  It also preserves binders: a variable node in $S$ is bound
either strictly above $m_1$, where nothing moves, or strictly below
$m_{2p}$, i.e.\ inside $S$, where $\varrho$ relocates binder and
occurrence together (Definition~\ref{def:endpath}(3); this also covers
$s$ itself).  Variable nodes outside the subtree at $m_1$ have their
binders outside it as well, and $w_1,\dots,w_k$ are untouched.

Two facts about $U$ are used repeatedly.  (U1) If $v\notin U$ is a
variable node, then no successor of $v$ is in $U$: the only node outside
$U$ with a successor in $U$ is $\lambda_0$, and $\lambda_0$ is a lambda
node.  The same holds for constant nodes and for the root $\varepsilon$.
(U2) If $v\notin U$ is a variable node, its binder is not in $\Gamma$:
for $v\in S$ by Definition~\ref{def:endpath}(3), and for $v$ outside the
subtree at $m_1$ because its binder lies above $v$, hence outside that
subtree, while $\Gamma\subseteq U$.

(1) Apply Definition~\ref{def:redundant} to $\iv1l$, and again to the
interval $\iv{j}{l}$ that clause~(c) produces, and so on; the left end
strictly increases, so this stops.  The $r$-th application gives the
least $a_r$ with $n_{a_r}=m_1$ after the previous $j_{r-1}$, and $j_r$
with a chain along $\gamma$ from $\pi_{a_r}$ to $\pi_{j_r}$ and
$\iv{a_r}{j_r}$ not contributing.  As $t$ solves $E$, the play ends at a
constant node with $[\exists]$, while $n_{j_r}=m_{2p}$ is a lambda node;
so $j_r<l$ and non-contribution means $r_{a_r}=r_{j_r}$.  The next
position after $\pi_{j_r}$ is at $s\ne m_1$, so $\sigma_{r+1}$ is
non-empty, and $a_1>1$ since $m_1\ne\varepsilon$.  By
Lemma~\ref{lem:states}, $r_{a_r}=r_{j_r}$ forbids $f$-moves in
$D_r$.  By the choice of the $a_r$ as least indices, no outer position
is at $m_1$.

Two more facts about the ends of $D_r$.  Since $n_{a_r}=m_1$ is not a
lambda node, Lemma~\ref{lem:pred} gives that $\pi_{a_r-1}$ is at
$\lambda_0$, with $\theta_{a_r-1}=\theta_{a_r}$ and
$r_{a_r-1}=r_{a_r}$; and $\pi_{a_r-1}$ is outer.  Along the chain, the
tables increase (Proposition~\ref{prop:mono}(2)).  By
Proposition~\ref{prop:tables}(1), $\theta_{a_r}$ is defined exactly on
the binding nodes strictly above $m_1$ and $\theta_{j_r}$ on those
above $m_{2p}$.  So
\[
\theta_{j_r}=\theta_{a_r}\cup E_r ,\qquad
\operatorname{dom}E_r\subseteq\Gamma. \tag{$*$}
\]

(2) Induction over the outer positions in order.  $\pi_1$ is at
$\varepsilon\notin U$ with the empty table.  Let $\pi_{i+1}$ be outer and
assume the claim for all outer positions before it.

\emph{Case $\pi_i$ is outer.}  Look at the move from $\pi_i$.  For the
$@$-, $\lambda$- and $f$-moves, $n_{i+1}$ is a successor of
$n_i\notin U$ and is not $m_1$ (as $\pi_{i+1}$ is outer); so by (U1),
or because the only successor of $\lambda_0$ is $m_1$, $n_{i+1}\notin
U$.  The table is unchanged, or, for the $@$-move, gets one entry with
inner table $\theta_1$, and $\pi_1$ is outer.  For a jump at a variable
node $n_i$ with binder $b$: by (U2), $b\notin\Gamma$, so by cleanness
$\theta_i(b)=((\suc v1,\dots),\theta_h)$ with $\pi_h$ outer.  Then
$n_{i+1}=\suc vg$ for some $g$, and $n_h=v$ by
Proposition~\ref{prop:tables}(4) (or $v=\varepsilon$, $h=1$).  As $\pi_h$
is outer, $v\notin U$, so $n_{i+1}\notin U$ by (U1).  The new table is
$\theta_h$, which is clean by the induction hypothesis, possibly extended
by an entry for $\suc vg\notin\Gamma$ with inner table $\theta_i$, where
$\pi_i$ is outer.  So it is clean.

\emph{Case $\pi_i$ is not outer.}  Then $\pi_i=\pi_{j_r}$ ends some
$D_r$, and $\pi_{i+1}$ is at $s\notin U$ with table $\theta_{j_r}$.  By
$(*)$ and $\theta_{a_r}=\theta_{a_r-1}$ it differs from the clean table
$\theta_{a_r-1}$ only by entries for nodes in $\Gamma$, which cleanness
ignores.  So it is clean.

(3) $\theta^\flat$ is well defined for clean $\theta$: the inner tables
$\theta_i$ are tables of outer positions, so clean by (2), and they lie
in lower layers.  The domain avoids $U$: by Proposition~\ref{prop:tables}(1)
it consists of binding nodes above an outer node $n\notin U$, and those
in $U$ are in $\Gamma$.  Arguments $\suc v1,\dots$ of entries avoid $U$
by (U1), since $v$ is the node of an outer position.  So $\varrho^{-1}$
applies everywhere.  Also $\varphi(\pi_1)=\varepsilon[u]\varnothing$.

Let $\pi_i,\pi_{i'}$ be consecutive outer positions.  We show that
$\varphi(\pi_{i'})$ is the next position of $\varphi(\pi_i)$ in
$\Gm(t'',E)$, for the same choice of $\forall$ if $n_i$ is a constant.

\emph{$i'=i+1$.}  The move from $\pi_i$ uses the label, successors and
binder of $n_i$, which $\varrho$ preserves, the state, which $\varphi$
keeps, and, for a jump, the entry of the binder $b\notin\Gamma$ (by (U2)).
That entry is $((\suc v1,\dots),\theta_h)$ in $\theta_i$ and
$((\varrho^{-1}\suc v1,\dots),\theta_h^\flat)$ in $\theta_i^\flat$.  The
new table is $\theta_h$ or its extension by
$((\suc{n_i}1,\dots),\theta_i)$ for $\suc vg$; in $\Gm(t'',E)$ it is
$\theta_h^\flat$ or its extension by
$((\varrho^{-1}\suc{n_i}1,\dots),\theta_i^\flat)$ for
$\varrho^{-1}\suc vg$.  This is $\theta_{i+1}^\flat$.

\emph{$i'=j_r+1$, $i=a_r-1$} (a deleted $D_r$).  $\pi_i$ is at
$\lambda_0$, so in $\Gm(t'',E)$ the $\lambda$-move leads from
$\varphi(\pi_i)$ to $\suc{\lambda_0}1[r_i]\theta_i^\flat$, and
$\suc{\lambda_0}1=m_1=\varrho^{-1}(s)$.  On the other side,
$\pi_{j_r+1}=s[r_{j_r}]\theta_{j_r}$ (a $\lambda$-move from
$m_{2p}$).  Here $r_{j_r}=r_{a_r}=r_i$ by (1), and
$\theta_{j_r}^\flat=\theta_{a_r}^\flat=\theta_i^\flat$ by $(*)$, because
$\flat$ deletes the entries for $\Gamma$.

So $\varphi(\pi)$ is a sequence of positions of $\Gm(t'',E)$ in which
each one is a next position of the previous one.  The last position of
$\pi$ is outer, so $\varphi(\pi)$ ends with the final state of $\pi$, and
earlier states are non-final.  It is a play.  By (1), every $f$-move
of $\pi$ is at an outer position, and $\varphi$ keeps it with the same
choice.

(4) Moves are deterministic except for $\forall$'s choices at $f$-moves.
Let $\pi''\in\Gm(t'',E)$ have choices $c_1,\dots,c_e$.  Let $\pi$ be the
play of $\Gm(t,E)$ in which $\forall$ makes choice $c_d$ at her $d$-th
$f$-move, for $d\le e$.  Plays are finite (Theorem~\ref{thm:game}), so
$\pi$ exists, provided each $c_d$ is admissible.  By (3), $\varphi(\pi)$
and $\pi''$ are plays of the same game that are built by the same moves
up to the first choice where they differ.  They make the same choices,
so they agree up to the end of the shorter one.  That end is a final
position, which has no successor, so $\varphi(\pi)=\pi''$.  Admissibility
holds by the same induction: the $d$-th $f$-move of $\pi$ is at an outer
position whose image is the $d$-th $f$-move of $\pi''$, with the same
state.
\end{proof}

\begin{remark}
The lemma does not use the order bound; at order $5$,
Lemma~\ref{lem:endpath5} makes it concrete.  For a $z$-path, $\Gamma$ is
empty, since $m_2$ is a dummy, and $\theta_{j_r}=\theta_{a_r}$.  For an
$x$-path, $\Gamma\subseteq\{m_2\}$, and $E_r$ is the entry for
$m_2=\suc nh$ created by the $y$-jump into $m_2$ inside $D_r$.  The
invariant~(2) is exactly what was missing in the paper's argument.
Tables of later positions may still contain entries created inside
$D_r$, but only for $m_2$, and no outer position consults them, because
every variable bound by $m_2$ lies in $U$.
\end{remark}

\begin{proposition}[S, Prop.~31]\label{prop:T2}
\Fixed{} If $t$ solves $E$ and $t\,T_2\,t''$, then $t''$ solves $E$.
\end{proposition}
\begin{proof}
By Theorem~\ref{thm:game}, $\forall$ loses $\Gm(t,E)$, so every play
ends with $[\exists]$.  By Lemma~\ref{lem:T2sim}(4) and~(3), every play
of $\Gm(t'',E)$ is $\varphi(\pi)$ for some $\pi\in\Gm(t,E)$, and it has
the same final state.  So $\forall$ loses $\Gm(t'',E)$, and $t''$ solves
$E$ by Theorem~\ref{thm:game}.
\end{proof}

\begin{definition}[Terse solution]\label{def:terse}
\New{} A solution $t\in T_A(\CE)$ of $E$ is \emph{terse} if neither $T_1$
nor $T_2$ applies to it.
\end{definition}

\begin{lemma}[Terse solutions exist]\label{lem:terse}
\New{} Starting from any solution in $T_A(\CE)$, every sequence
of $T_1$- and $T_2$-steps is finite, and its last term is a terse
solution.  So if $E$ has a solution, it has a terse one.
\end{lemma}
\begin{proof}
Each step decreases $\mu$ (Proposition~\ref{prop:T1} and the definition
of $T_2$), and $\mu$ is well founded, being a lexicographic triple of
natural numbers.  Each step preserves solutions
(Propositions~\ref{prop:T1} and~\ref{prop:T2}).  The last term admits no
step, so it is terse.
\end{proof}

\begin{remark}
The rest of the proof uses only terseness, never minimality.  Where the
paper argues ``$T_2$ would give a smaller solution'', it is enough that
$T_2$ does not apply.  That $T_2$ preserves solutions
(Proposition~\ref{prop:T2}) is used only in Lemma~\ref{lem:terse} and
Step~4 of Section~\ref{sec:architecture}.
\end{remark}

\begin{proposition}[Strips of $z$-nodes; S, Prop.~32]\label{prop:stacking}
\Fixed{} Let $t$ be a terse solution, and let $m_1,\dots,m_p$ be
$z$-nodes of $t$ carrying the same variable with the same binder $b$,
with $m_{i+1}$ below $\suc{m_i}{l_i}$ for $i<p$.
\begin{enumerate}
\item If $l_1=\dots=l_{p-1}$, then $p\le\depth(u)+1$.
\item In general, $p\le q\,\depth(u)+1$.
\end{enumerate}
\end{proposition}
\begin{proof}
All $m_i$ lie on one path, so $m_p$ lies below $\suc{m_i}{l_i}$ for every
$i<p$.  Let $e$ be an index that occurs most often among
$l_1,\dots,l_{p-1}$, and let $I=\{i<p\mid l_i=e\}$.  A $z$-variable has
at most $q$ arguments, so $|I|\ge(p-1)/q$, and $|I|=p-1$ in case~(1).
We show $|I|\le\depth(u)$.

\emph{A witness at $m_p$.}  $(m_p,\suc{m_p}e)$ is an end path
(Lemma~\ref{lem:endpath5}).  $t$ is terse, so it is not redundant.
Unfolding Definition~\ref{def:redundant}, some play $\pi$ has a position
$\pi_j$ at $m_p$ at which (a) or (b) fails.  $m_p$ is end, so $\pi_j$
has at most one later child (Fact~\ref{fact:25}).  Hence the excursion
at $\pi_j$ (Definition~\ref{def:excursion}) does not return to
$\suc{m_p}e$, or it returns there and contributes.  (By
Lemma~\ref{lem:pred}, $r_j=r_{j+1}$, and the state at the child equals
the state just before it.)

\emph{The earlier visits.}  $\pi_{j-1}$ is at the lambda node above
$m_p$ (Lemma~\ref{lem:pred}).  By Lemma~\ref{lem:ancestors} there is a
chain for $\pi_{j-1}$ starting at a position at the $x$-node above $b$.
For $i\in I$ it contains a position $\pi_{a_i}$ at $m_i$, followed by
its child $\pi_{c_i}$ at $\suc{m_i}e$.  Along a chain the indices
increase and the tables extend (Proposition~\ref{prop:mono}), so
$a_i<c_i\le a_{i'}$ for $i<i'$ in $I$, and the position of the chain at
$b$ has the same entry for $b$ as $\pi_{a_i}$ and $\pi_j$.  So it is the
common parent of $\pi_{a_i}$ and $\pi_j$.  By Fact~\ref{fact:25},
$\pi_{c_i}$ is the only later child of $\pi_{a_i}$, so the excursion at
$\pi_{a_i}$ is $\iv{a_i+1}{c_i-1}$ and returns to $\suc{m_i}e$.

\emph{Each of these excursions contributes.}  Fix $i\in I$ and apply
Proposition~\ref{prop:24} to $\pi_{a_i}$ and $\pi_j$.  It gives $k$ such
that $\iv{a_i+1}{a_i+k}$ and $\iv{j+1}{j+k}$ correspond.  By its proof,
the corresponding segments stop at the latest at the first return, so
$a_i+k<c_i$.
\begin{itemize}
\item Case~(1): $\pi_{a_i}$ is the parent of $\pi_{a_i+k+1}$, so
  $a_i+k+1=c_i$, and $\pi_j$ returns at $\pi_{j+k+1}$ to $\suc{m_p}e$.
  By the witness property, $\iv{j+1}{j+k}$ contributes, so it contains
  a position at a constant of positive arity (Lemma~\ref{lem:states}).
  The corresponding position in $\iv{a_i+1}{a_i+k}$ is at the same
  node, so the excursion at $\pi_{a_i}$ contributes.
\item Case~(2): $n_{a_i+k}$ is a constant node.  It cannot have arity
  $0$: the $a$-move would end the play before $\pi_j$.  So $\pi_{a_i+k}$
  makes an $f$-move to a successor of a constant.  That position is not
  $\pi_{c_i}$, which is entered by a $v$-jump.  So the excursion at
  $\pi_{a_i}$ contains the $f$-move and contributes
  (Lemma~\ref{lem:states}).
\end{itemize}

\emph{Counting.}  The intervals $\iv{a_i+1}{c_i-1}$, $i\in I$, are
disjoint and ordered, and each has different states at its ends.
Lemma~\ref{lem:disjoint} gives $|I|\le\depth(u)$.  Hence
$p-1\le q\,\depth(u)$, and $p-1\le\depth(u)$ in case~(1).
\end{proof}

\begin{remark}
The paper claims $p\le\depth(u)$.  The argument gives one contributing
excursion for each $m_i$ with $i<p$ and the right successor index, not
for $m_p$ itself.  The factor $q$ in~(2) comes from the choice of the
index $e$: when the witness at $m_p$ returns to some $\suc{m_p}{e'}$
with $e'\ne e$, Proposition~\ref{prop:24} says nothing about the
$m_i$ with $l_i=e'$.
\end{remark}

%======================================================================
\section{Partitioning play [S, \S6.3]}\label{sec:partition}
%======================================================================

Throughout this section $t$ solves $E$ and $\pi\in\Gm(t,E)$ has length
$l$.

\begin{definition}[Stages; S, Def.~43]\label{def:stages}
\Fixed{} Stage~$1$ is $\iv12$.  Let stage $N$ end at $\pi_{i-1}$, a
lambda node of $t$, so $n_i$ is its successor.
\begin{enumerate}
\item If $n_i$ is a constant node, stage $N+1$ is $\iv i{i+1}$.
\item If $n_i$ is a variable node, $n_{i+1}$ is a lambda node of some
  $w_q$ (Lemma~\ref{lem:alt}).  Let $\pi_{j_1},\dots,\pi_{j_{2r}}$,
  $j_1=i+2$, be the greedy recurrent chain
  (Proposition~\ref{prop:29}).  Let $p$ be the least $p'\in\{1,\dots,r\}$
  such that $\pi_{j_{2p'}+2}$ is a child of $\pi_i$, or $p=r$ if there is
  none.  Stage $N+1$ is $\iv i{i+3}$ if $p=0$, and $\iv i{j_{2p}+2}$
  otherwise.  The chain $\pi_{j_1},\dots,\pi_{j_{2p}}$ is
  \emph{associated} with the stage.
\end{enumerate}
\end{definition}

\begin{remark}
The paper's definition requires that ``$\iv{j_1}{j_{2p}}$ does not
contain a child of $\pi_i$''.  The text after it uses the cut ``first
$p$ with $\pi_{j_{2p}+2}$ a child of $\pi_i$'', which we adopt.  The
first member $\pi_{j_1}=\pi_{i+2}$ is at a variable or constant node of
$w_q$, not at a lambda node (the margin note ``$\pi_{j_1}$ jest
lambd\k{a}'' on p.~46 of the annotated copy does not fit).  This follows
from Definition~\ref{def:chain}: odd members of a chain are parents of
lambda positions.
\end{remark}

\begin{fact}[S, Fact~33]\label{fact:33}
\Paper{} The stages are consecutive, non-overlapping intervals covering
$\pi$.  Every non-final stage ends at a lambda node of $t$, and the
partition is unique.
\end{fact}
\begin{proof}
Assume stage $N$ ends at a lambda node of $t$.
\begin{itemize}
\item \emph{Constant case:} the stage ends at a successor of the
  constant, or is final.
\item \emph{Variable case, $p=0$:} then $r=0$.  The node $n_{i+2}$ cannot
  be a constant of positive arity, since that would give a recurrent
  chain of length~$2$.  So $\pi_{i+3}$ is a jump into $t$, or the play
  ends at $a$.
\item \emph{Variable case, $p\ge1$:} either $\pi_{j_{2p}+2}$ is a child of
  $\pi_i$, which lies at a successor of $n_i$ in $t$, or $p=r$.  In the
  latter case $\pi_{j_{2r}+1}$ is a variable node, since a constant would
  extend the chain, so $\pi_{j_{2r}+2}$ jumps into $t$, or the play ends.
\end{itemize}
Uniqueness follows from the uniqueness of the greedy chain.
\end{proof}

\begin{definition}[S, Def.~44]
A stage $\iv ij$ is \emph{final} if $r_j$ is final.  A non-final stage
\emph{contributes} if $\iv ij$ does, and a final one if $\iv i{j-1}$
does.  Equivalently (Lemma~\ref{lem:states}), a stage contributes iff
it contains a position at a constant node of positive arity.
\end{definition}

\begin{corollary}\label{cor:contrib-stages}
\New{} In every play at most $\depth(u)$ stages contribute.  Over the
tree of all plays at most $|u|-\br(u)$ distinct stages contribute.
\end{corollary}
\begin{proof}
Lemmas~\ref{lem:disjoint} and~\ref{lem:global}.
\end{proof}

\begin{proposition}[S, Prop.~34]\label{prop:34}
\Paper{} Let $\iv ij$ be stage $N$.
\begin{enumerate}
\item If $n_i$ is labelled $a:\Ground$, then $N$ is final.
\item If $N$ is not final and $n_i$ is a constant or $x$-node,
  then $\pi_j$ is a child of $\pi_i$.
\item If $N$ does not contribute and $n_i$ is embedded, then $\pi_j$ is a
  child of $\pi_i$.
\item If $N$ is not final and $n_i$ is a variable node, then $n_j$ is a
  successor of a node in the same family as $n_i$.
\end{enumerate}
\end{proposition}
\begin{proof}
(1) The $a$-move ends the play.

(2) The constant case is Definition~\ref{def:parent}(2).  In the
$x$-case, the entry for the root of $t$ is fixed, so $n_{i+1}$ is
the root of $w_q$.  The constants met on the chain bind nothing.  So the
first variable met is bound by that root, and its jump is to a child of
$\pi_i$; this is the cut in Definition~\ref{def:stages}.

(3) By Fact~\ref{fact:27} the excursion is complementary to an earlier
one at the embedding node, which returned to a child.

(4) Extend the associated chain upward to the root of $w_q$ and use
Proposition~\ref{prop:28}.
\end{proof}

\begin{lemma}[Stages at order $5$]\label{lem:stages5}
\New{} Every stage other than stage~$1$ is a \emph{constant stage}
$\iv i{i+1}$, an \emph{$x$-stage} or a \emph{$z$-stage}, according to the
label of $n_i$.  Let $N=\iv ij$ be non-final.
\begin{enumerate}
\item If $n_i$ is an $x$-node, then $n_j$ is a successor $\lambda\vec z$
  of $n_i$, entered by a $y$-jump, and $\pi_j$ is a child of $\pi_i$.
  $N$ does not contribute iff the body $H$ of the $w$ entered at
  $\pi_{i+1}$ has a $y$-variable $y_h$ as head.  Then $N=\iv i{i+3}$ and
  $n_j=\suc{n_i}h$.
\item If $n_i$ is a $z$-node of the family $F(n_0)$, then $n_j$ is
  either a dummy lambda below a $z$-node of $F(n_0)$, entered by a
  $v$-jump, or a successor $\lambda\vec z$ of $n_0$, entered by a
  $y$-jump (a \emph{re-entry}).  If $N$ does not contribute and $n_i$ is
  embedded, it is the dummy lambda $\suc{n_i}e$ below $n_i$ itself, a
  child of $\pi_i$.
\end{enumerate}
\end{lemma}
\begin{proof}
By Fact~\ref{fact:33} a non-final stage ends at a lambda node of $t$.
By Corollary~\ref{cor:jumps}, lambda nodes of $t$ are entered from $w$ by
$y$-jumps (at successors of $x$-nodes) or $v$-jumps (at dummies below
$z$-nodes).

(1) Proposition~\ref{prop:34}(2).  $\pi_{i+2}$ is at the head of $H$
(the $w$ has no free variable, so the head is a $y$-variable or a
constant).  A constant head of positive arity contributes, and one of
arity~$0$ ends the play.  A head $y_h$ jumps at $\pi_{i+3}$ to
$\suc{n_i}h$, which is a child of $\pi_i$, and this is the cut.

(2) Proposition~\ref{prop:34}(4): $n_j$ is a successor of a node of
$F(n_0)$, and the only $x$-node of $F(n_0)$ is $n_0$
(Remark on families in Section~\ref{sec:uniformity}).  The last claim is
Proposition~\ref{prop:34}(3).
\end{proof}

\begin{definition}[Comb, path; S, Def.~45]
Let $\iv ij$ be stage $N$.
\begin{itemize}
\item The \emph{comb} $\comb(N)$ is $n_i$ with its successors.
\item The \emph{path} $\pth(N)$ is empty for constant and $@$ stages.
  For a variable stage with associated chain
  $\pi_{j_1},\dots,\pi_{j_{2p}}$ it is
  $n_{i+1},n_{j_1},\dots,n_{j_{2p}},n_{j-1}$, a downward path in some
  $w_q$.
\end{itemize}
Stages inherit \emph{constant/variable}, \emph{level},
\emph{embedded} and \emph{family} from the root $n_i$ of their comb.
\end{definition}

\begin{definition}[Partition tree; S, Def.~46]\label{def:ptree}
The nodes are the stages of $\pi$.  For each non-final stage $N=\iv ij$
there is an edge $N+1\ra\pi K$, where $K\le N$ is the largest stage
whose comb contains $n_j$ and such that $N\ra\pi^*M$ and $K\ra\pi^*M$
for some $M\le N$.
\end{definition}

\begin{remark}[Degeneracy]\label{rem:ptree}
\Fixed{} Every stage reaches stage~$1$, so the side condition always
holds with $M=1$.  As written, the edge goes to ``the latest stage whose
comb contains $n_j$''.  The same node can be the comb root of several
stages, visited with different tables, so this is not the intended
relation.  We use the following reading.
\end{remark}

\begin{definition}[Partition tree, proposed reading]\label{def:ptree2}
\New{} For a non-final stage $N=\iv ij$, let $\pi_h$ be the parent of
$\pi_j$.  Put $N+1\ra\pi K$, where $K$ is the stage containing $\pi_h$.
Gluing $\comb(N+1)$ below the lambda leaf $n_j$ of $\comb(K)$ along every
edge is intended to rebuild the part of $t$ visited by $\pi$.
\end{definition}

\renewcommand{\gapname}{[Gap-Partition]}
\begin{gapenv}[Partition tree]\phantomsection\label{gap:ptree}
\Open{} Needed:
\begin{enumerate}[(i)]
\item $\pi_h$ is the first position of $K$, so $n_j\in\comb(K)$;
\item the proof of Theorem~\ref{thm:small} goes through with
  Definition~\ref{def:ptree2};
\item every visited node of $t$ lies in the comb of some stage.
\end{enumerate}
\end{gapenv}

\begin{definition}[Hyper-partition tree]\label{def:hyper}
\New{} Stages $N$ of $\pi$ and $N'$ of $\pi'$ are \emph{identical}
($N\equiv N'$) if $N=N'$ as numbers and the first $N$ stages of $\pi$
and of $\pi'$ are the same intervals of the same positions.  The
\emph{hyper-partition tree} $\mathcal H$ has as nodes the
$\equiv$-classes of stages of all plays of $\Gm(t,E)$, i.e.\ it is the
trie of the stage sequences of the plays.  It carries two relations:
\begin{itemize}
\item the \emph{play order}: $[N+1]$ is a successor of $[N]$;
\item the \emph{partition edges} $[N+1]\ra{}[K]$ of
  Definition~\ref{def:ptree2}, taken from any play through $[N+1]$.
\end{itemize}
Two nodes are \emph{comb-equivalent} ($[N]\sim[N']$) if they are
successors of the same node and start at the same position.  A
\emph{separator} is a $\sim$-class with more than one node.
\end{definition}

\begin{remark}
One cannot cut the common prefix of two plays into stages once and
for all.  A stage can run past the point where the plays diverge, and
then it ends differently in each.  In Example~\ref{ex:obstruction},
plays A and B agree up to $\pi_{11}$; $S_4$ is $\iv9{18}$ in~A and
$\iv9{12}$ in~B.  These are two nodes of $\mathcal H$, comb-equivalent,
forming a separator.  Proposition~\ref{prop:hyper}(4) shows that this
is the only way stages of plays with a common history can differ.
\end{remark}

\begin{proposition}[Properties of $\mathcal H$]\label{prop:hyper}
\New{}
\begin{enumerate}
\item The play order makes $\mathcal H$ a tree rooted at stage~$1$.
  Each play is a branch of it.
\item The partition edges are well defined: the edge out of $[N+1]$
  does not depend on the play chosen, and it points to an ancestor of
  $[N+1]$ in the play order.  So they also form a tree, which
  restricted to one play is that play's partition tree.
\item Comb-equivalent nodes have the same comb.
\item Every separator contains a node that contributes.
\end{enumerate}
\end{proposition}
\begin{proof}
(1) $\equiv$ is defined by prefixes of stage sequences.

(2) The edge out of $N+1$ is determined by the last position $\pi_j$ of
$N$ and its parent $\pi_h$, with $h\le j$
(Definition~\ref{def:ptree2}).  Both lie in the first $N$ stages,
which are the same in all plays through $[N+1]$.  The target $K$ is
the stage containing $\pi_h$, so $K\le N$.

(3) The comb is determined by the first position.

(4) Let $[N_A]\sim[N_B]$ be different, coming from plays $A$
and $B$, with common first position $\pi_i$ and equal histories up to
it.  Play is deterministic except for $\forall$'s choice at an
$f$-move (Definition~\ref{def:moves}).  Let $\pi_{c-1}$ be the first
position at or after $\pi_i$ at which $\forall$ chooses differently in $A$
and $B$.  Say $N_B$ is the longer of the two stages.  Suppose $N_B$
ends before $\pi_{c-1}$.  Then its greedy chain, and the positions that
decide its cut (Definition~\ref{def:stages}), lie in the common prefix.
The greedy construction in $A$ would therefore follow the same chain
and make the same cut, so $N_A=N_B$, a contradiction.  Hence $N_B$
reaches $\pi_{c-1}$.  An $f$-position is neither final nor at a lambda
node, so it is not the last position of $N_B$.  So $N_B$ contains the
$f$-move at $\pi_{c-1}$ and contributes ([S, Def.~44]).

The step ``lie in the common prefix'' holds because everything that
decides the stage is local to it.  Proposition~\ref{prop:29} looks for
the \emph{first} later child position.  The recurrence property
(Definition~\ref{def:recurrence}) compares the interval
$\iv{m+1}{m+k}$ with an earlier one, where $j_{2i}=m+k+1$.  The clauses
of complementarity (Proposition~\ref{prop:24}) mention positions only
up to $\pi_{m+k+1}$.  All of these lie at or before the end of the
stage.
\end{proof}

\begin{remark}
With Definition~\ref{def:47a}, unfoldability is a property of a
$\sim$-class: it requires that no node of the class contributes.  By
Proposition~\ref{prop:hyper}(4), separators are never unfolded.  This is
what the construction in Section~\ref{sec:thm37} requires (``separators
must be kept'').  Their number is not bounded here; the proof bounds
strips instead (Section~\ref{sec:architecture}).
\end{remark}

\begin{fact}[S, Fact~35]\label{fact:35}
\Sketch{} If stage $N=\iv ij$ contributes, then one of the following
holds:
\begin{enumerate}
\item $N$ is a constant stage;
\item $\pth(N)$ contains a constant node;
\item some $n_k$, $i<k<j$, is a constant node of $t$ in the comb of an
  earlier stage of the same family.
\end{enumerate}
The third case uses \GapPartition{}(iii).
\end{fact}

%======================================================================
\section{Unfolding [S, \S6.4]}\label{sec:unfolding}
%======================================================================

\subsection{The paper's definition}

\begin{definition}[Unfolding; S, Def.~47]\label{def:unfold}
Let $N_1<\dots<N_k$ be stages of $\pi$ with combs $c_i$ and paths
$\gamma_i$, such that
\begin{enumerate}[(a)]
\item they belong to the same family and do not contribute;
\item $c_i\ne c_j$ and $\gamma_i\ne\gamma_j$ for $i<j$;
\item for every $j\le k$ (except $j=k$ if $N_k$ is final) there is
  $i\le j$ with $N_j+1\ra\pi N_i$;
\item for every $1<j\le k$ there is $i<j$ such that the root of
  $\gamma_j$ is a successor of the last node of $\gamma_i$.
\end{enumerate}
The \emph{enlarged comb} is built as follows: $(c_1)=c_1$, and
$(c_1\dots c_{j+1})$ places a fresh copy of $c_{j+1}$ below the copy of
the leaf at which $N_j$ finishes.  Variables of the copy are bound by
the copies of their binders when these lie in some $c_i$.  The
\emph{enlarged tree} $(\gamma_1\dots\gamma_j)$ is the union of the paths.
The new stage $(N_1\dots N_j)$ starts at the root of $(c_1\dots c_j)$
and finishes at one of its leaves.  Edges into $N_i$ given by (c) are
redirected to the new stage.
\end{definition}

\begin{fact}[S, Fact~36]\label{fact:36}
\Open{} For each $j\le k$, the play of $(N_1\dots N_j)$ does not
contribute, and it finishes at node $n$ of $c_i$ in $(c_1\dots c_j)$ iff
$N_j$ finishes at $n$ of $c_i$.
\end{fact}

\begin{remark}
The paper does not define the game of the unfolded term (``we do not
provide the game to save notation'').  It also does not prove that the
unfolded term solves $E$.  Section~\ref{sec:nonlinear} shows that, as
written, Definition~\ref{def:unfold} can produce a non-term.
\end{remark}

\begin{definition}[c-embedded; S, Def.~48]
$(c_1\dots c_k)$ is \emph{c-embedded} if its only free variable
occurrence is the head of $c_1$ and some $\alpha$-equivalent enlarged
comb $(c'_1\dots c'_k)$ lies above it in the partition tree.
\end{definition}

\begin{lemma}[Shape of c-embedded unfoldings]\label{lem:cembshape}
\New{} Let $(c_1\dots c_k)$ be c-embedded, $k\ge2$.  Then $c_1$ is the
comb of an $x$-stage at a node $n$, and $c_2,\dots,c_k$ are combs of
$z$-stages whose heads are bound by one successor $\suc nh$ of $n$, and
which lie one below another on one path.  If no $N_j$ is embedded
because of an earlier $N_i$, their variables are pairwise distinct, so
$k\le q+1$, and the enlarged comb has depth at most $2(q+1)$.
\end{lemma}
\begin{proof}
A $z$-comb binds nothing (Corollary~\ref{cor:binders}), and a constant
comb of a non-contributing stage binds nothing either.  The head of
$c_j$, $j\ge2$, is bound in some earlier $c_i$, so in the only comb that
binds anything, $c_1$, which must therefore be an $x$-comb.  Each $c_{j+1}$
is placed below the leaf at which $N_j$ finishes, so all $c_j$, $j\ge2$,
lie below one successor $\suc nh$ and are bound by it.  It binds at most
$q$ variables.  Each comb contributes a variable node and a lambda node
to the depth.
\end{proof}

\subsection{Linear heads}

We write $t=\lambda x_1\dots x_k.\,S$ and $w_i=\lambda\vec y.\,H_i$.

\begin{definition}[Linear and non-linear heads]\label{def:linear}
\New{} The head $y_j$ of $w_i=\lambda\vec y.\,y_j B_1\dots B_p$ is
\emph{linear} if $y_j$ does not occur in $B_1,\dots,B_p$, and
\emph{non-linear} otherwise.  Repetitions of other variables $y_h$,
$h\ne j$, do not matter here.
\end{definition}

\begin{definition}[One-step unfolding]\label{def:onestep}
\New{} Let $n=x_i\,A_1\dots A_m$ be an $x_i$-node of $t$ with
$A_j=\lambda\vec z.\,S_j$, and let $m=z_l\,R_1\dots R_s$ be a node in
$S_j$ with $z_l$ bound at $\suc nj$.  The \emph{unfolding of $t$ at
$m$} is
\[
t'=t\bigl[m:=x_i\,A_1\dots A_{j-1}\,(\lambda\vec z'.\,z'_l\,\vec R')\,
A_{j+1}\dots A_m\bigr],
\]
where $\vec z'$ is fresh and $\vec R'$ is $\vec R$ with the occurrences of
$\vec z$ bound at $\suc nj$ renamed to $\vec z'$.  So the new $x_i$-node
at $m$ carries:
\begin{itemize}
\item under its $h$-th lambda, $h\ne j$: a copy of the whole argument
  $A_h$;
\item under its $j$-th lambda: the node $z'_l$, with the original
  arguments $\vec R'$ under its dummy lambdas.
\end{itemize}
The node $n$ and its arguments stay unchanged.  This is
Definition~\ref{def:unfold} with $k=2$; for $k>2$ the operation is
repeated inside the copy.
\end{definition}

\begin{remark}[The comb view]
In terms of combs, the copy of $\comb(N_1)$ (the $x_i$-node with its
lambdas) is placed at $m$, and the copy of $\comb(N_2)$ (the $z_l$-node
with its dummy lambdas) under its $j$-th lambda.  Under the other
lambdas hang the stages whose partition edges lead there.  A lambda that
no play reaches gets nothing in the comb view; in
Definition~\ref{def:onestep} it carries a copy of the original
subterm, which is unvisited and is replaced by $b$ once $T_1$ is
applied.  So the comb view describes exactly the visited part of $t'$.
\end{remark}

\begin{definition}[Unfolding regime]\label{def:regime}
\New{} Order the nodes of a tree by their addresses: $m\prec m'$ if the
address of $m$ is lexicographically larger than that of $m'$, where a
proper prefix counts as smaller.  So a node comes after all its
descendants and after everything to its right: $\prec$ goes from the
bottom right of the tree to the left and up.

A $z$-node $m$ is \emph{unfoldable} in $t$ if it is bound at $\suc nj$
for an $x$-node $n$ and there are stages $N_1$ at $n$ and $N_2$ at $m$
satisfying Definition~\ref{def:unfold}(a)--(d) with $k=2$, read over
$\mathcal H$ (Definition~\ref{def:47a}).

The \emph{regime} repeats: if $t$ has an unfoldable node, unfold $t$ at
the $\prec$-least one (Definition~\ref{def:onestep}), recompute the
game and $\mathcal H$ for the result, and continue; otherwise stop.
\end{definition}

\begin{remark}
The regime makes the construction deterministic: every step is fixed
by $t$.  Unfolding at $m$ changes only the subtree at $m$, so the part
of the tree that comes $\prec$-before $m$ outside that subtree is not
changed.  Whether the regime terminates is part of \GapCounting{}(i).
\end{remark}

\begin{theorem}[Linear unfolding is sound]\label{thm:linunfold}
\New{} \Coq{} If $y_j$ is a linear head, then $t\,\vec w$ and
$t'\,\vec w$ have a common $\beta$-reduct.  Hence $t'$ solves $E$ iff
$t$ does.
\end{theorem}
\begin{proof}
\begin{enumerate}
\item Contract $w_i\vec A$ at $n$: it reduces to
  $S_j\{\vec z:=\vec B\{\vec y:=\vec A\}\}$.  By linearity, $S_j$
  occurs once.
\item At $m$, the reduct of $t$ holds
  $B_l\{\vec y:=\vec A\}\,\vec R\{\vec z:=\vec B\{\vec A\}\}$.
\item The reduct of $t'$ holds $w_i\vec A''$, with $\vec A''$ equal to
  $\vec A$ off~$j$.  This reduces to the same term, because
  $\vec B\{\vec y:=\vec A''\}=\vec B\{\vec y:=\vec A\}$ by linearity.
\end{enumerate}
Confluence finishes the proof.  In Coq this is \texttt{unfold\_sound}
(\texttt{Unfolding.v}), proved on raw de~Bruijn terms using
\texttt{confluence} (\texttt{Confluence.v}), with no axioms.  The lnf
version is still to be done.
\end{proof}

\subsection{Non-linear heads: the obstruction and its repair}\label{sec:nonlinear}

\begin{example}[Definition~47 yields no term]\label{ex:obstruction}
\Fixed{} Let $E$ be $X\,w=f\,(f\,a\,d)\,d$ with
$w=\lambda y.\,y\,(\lambda v.\,y\,(\lambda v'.\,v))$ and
$t=\lambda x.\,x\,(\lambda z.\,f\,(z\,a)\,d)$.  The head $y$ is
non-linear.  Number the nodes:
\begin{itemize}
\item in $t$: (1)$\lambda x$, (2)$x$, (3)$\lambda z$, (4)$f$, (5)$\lambda$,
  (6)$z$, (7)$\lambda$, (8)$a$, (9)$\lambda$, (10)$d$;
\item in $w$: (11)$\lambda y$, (12)$y$, (13)$\lambda v$, (14)$y$,
  (15)$\lambda v'$, (16)$v$.
\end{itemize}
In play~B, $\forall$ chooses branch~$1$ at the root $f$ and branch~$2$
at the re-entry:
\[
(0)\,(1)\,(2)\,(11)\,(12)\,(3)\,(4)\,(5)\,(6)\,(13)\,(14)\,(3)\,(4)\,(9)\,(10)\,(10)[\exists].
\]
Its stages under Definition~\ref{def:stages} are
$S_1=\iv12$, $S_2=\iv36$ ($x$, ending at (3)), $S_3=\iv78$ ($f$),
$S_4=\iv9{12}$ ($z$; the chain from $\pi_{11}$ is empty, so it ends at
the re-entry into (3)), $S_5=\iv{13}{14}$ and $S_6=\iv{15}{16}$.

Unfolding $S_2$ into $S_4$ satisfies (a)--(d) of
Definition~\ref{def:unfold}.  In $(c_2c_4)$ the copy of (3) carries the
copy of the $z$-node.  The redirected edge $S_5\to S_2$ must hang the
comb (4) under the same copy of~(3).  So the result is not a term.
\end{example}

\begin{definition}[Tree-wide non-contribution]\label{def:47a}
\Fixed{} Work in the hyper-partition tree (Definition~\ref{def:hyper}).
In Definition~\ref{def:unfold}(a), read ``do not contribute'' as ``no
node of the $\sim$-class of the stage contributes'', i.e.\ the stage
does not contribute in \emph{any} play passing through its comb.  This is
also how Stirling's LMCS~2009 version works (Defs.~9.1, 9.5(3)).
\end{definition}

In Example~\ref{ex:obstruction}, plays~A (branch~1 at the re-entry) and
B coincide up to $\pi_{11}$, so they share $S_4$.  In~A the chain
$\pi_{11},\pi_{16}$ is recurrent: $\iv69$ and $\iv{12}{15}$ are
complementary, with common parent $\pi_4$.  So $S_4=\iv9{18}$ passes
the constant at~(4) and contributes.  Under Definition~\ref{def:47a},
$S_2$ is therefore not unfoldable.

\begin{proposition}[Re-entries contribute somewhere]\label{prop:reentry}
\New{} Use Definitions~\ref{def:stages}
and~\ref{def:47a}.  Let $N_1=\iv{i_1}{i_1+3}$ be an $x$-stage at $n$ with
head $y_j$ at $\pi_{s-1}$, ending at $\pi_s$ at $\suc nj$.  Let $N_2$
start at $\pi_i$ at a $z_l$-node, and suppose its path starts at a
successor of the head $y_j$ node (Definition~\ref{def:unfold}(d)).
Suppose that in some play~B, $N_2$ ends without contributing at a
re-entry $\pi_e=\pi_{y+1}$ into $\suc nj$ from an inner $y_j$-node at
$\pi_y$.  Assume also the \emph{same-visit} condition: $\pi_{s-1}$ and
$\pi_y$ have the same parent (the visit of $\lambda\vec y$ at
$\pi_{i_1+1}$).  Then $N_2$ contributes in some play~A that shares $N_2$
with~B.
\end{proposition}
\begin{proof}
\emph{Step 1.} By Proposition~\ref{prop:24} for $\pi_{s-1},\pi_y$, there
is $k$ with $\iv s{s-1+k}$ and $\iv{y+1}{y+k}$ corresponding, and
either (1) both runs return to children of $\pi_{s-1}$, resp.\ $\pi_y$,
with the same index (a $\vec z$-node jumps into $B_m$, resp.\ $C_m$), or
(2) they part at a constant~$c$.

\emph{Step 2: case (1) is impossible in B.}  In case~(1) the intervals
are complementary (Definition~\ref{def:compl}).  So the link
$(\pi_y,\pi_{y+k+1})$ is recurrent.  It is also the first child of
$\pi_y$, since successors of $n_y$ are reached only through $\vec z$.
No child of $\pi_i$ occurs in $\iv y{y+k+1}$, since such a child needs a
variable bound in $B_l$, and these occur only inside $\vec C$.  So the
greedy chain of $N_2$ extends past $\pi_y$, and $N_2$ does not end
at~$\pi_e$, a contradiction.

\emph{Step 3: case (2) gives A.}  ``Final'' at $c$ is impossible: the
earlier run would then be final at $\pi_{s-1+k}$ with $s-1+k\le i$, but
it went on to reach the node of $N_2$.  So $\forall$ chose a different
branch at $c$, and $c$ has positive arity.  Let A agree with B up to
$\pi_{y+k}$ and copy the earlier choices from then on; such a play exists
by Fact~\ref{fact:branches}.  By Proposition~\ref{prop:24} again, A
reaches case~(1), at the latest at the copy of the $z_l$-node itself.
So A's chain extends past $\pi_y$, and A's $N_2$ contains $c$.  It
therefore contributes (Lemma~\ref{lem:states}).  A and B coincide up to
$\pi_{y+k}$, beyond $\pi_i$, so they share $N_2$.
\end{proof}

\begin{corollary}[Cross-visit case]\label{cor:crossvisit}
\New{} Drop the same-visit condition: two stages end at the same comb
leaf $L=\suc{n_r}h$ as children of different visits $V_1,V_2$ of $n_r$.
Compare the runs from the two entries into $L$ by
Proposition~\ref{prop:24star}.
\begin{itemize}
\item A split at $L$'s own binder gives complementary intervals
  (Corollary~\ref{cor:24prime}), and the argument of
  Proposition~\ref{prop:reentry} goes through.
\item A $\forall$ split is handled as in Step~3.
\item The new case is a split at an outer binder.  The
  root $x$'s have a single visit, so the only outer binders that can
  differ are the $\lambda\vec z'$ of level-$2$ variables bound above
  $n_r$.
\end{itemize}
Hence, if no $\lambda\vec z'$ lies on the path from the root of $t$ to
$n_r$, the contexts of $V_1,V_2$ are $\equiv$ and the conclusion of
Proposition~\ref{prop:reentry} holds.
\end{corollary}

\renewcommand{\gapname}{[Gap-Unfolding]}
\begin{gapenv}\phantomsection\label{gap:unfold}
\Open{}
\begin{enumerate}[(i)]
\item A precise game for the unfolded term, and Fact~\ref{fact:36}.
\item Soundness of unfolding in general.  It is proved here for linear
  heads (Theorem~\ref{thm:linunfold}).  For non-linear heads, with
  Definition~\ref{def:47a}, the conflict of
  Example~\ref{ex:obstruction} is excluded in the same-visit case
  (Proposition~\ref{prop:reentry}).
\item The remaining cross-visit case of Corollary~\ref{cor:crossvisit}:
  $n_r$ lies below some $\lambda\vec z'$ that its two visits bind
  differently, and play reaches that $\vec z'$ before $\forall$ splits.
\item Consistency across plays: Definition~\ref{def:unfold} composes
  combs per node, not per visit.
\item Compression.  Tree-shaped unfolding (Definition~\ref{def:treeunf})
  is sound under silence (Theorem~\ref{thm:treeunf}); what is open is
  the cut that turns it into Definition~\ref{def:onestep}
  (Remark~\ref{rem:compress}).
\end{enumerate}
\end{gapenv}

\begin{remark}[Alternatives considered]
\begin{itemize}
\item \emph{Def.~43$'$.} A stage absorbs re-entries into already visited
  $\lambda$-nodes and ends only at (E1) a child of $\pi_i$, (E2) a
  first entry into a new $\lambda$-node, (E3) a $\forall$ split inside
  an absorbed interval, or (F) a final state.  This follows LMCS
  Def.~8.13 (\texttt{order5-unfolding.md}~\S9).  It also blocks
  Example~\ref{ex:obstruction}, but it requires re-proving
  Fact~\ref{fact:33}, Proposition~\ref{prop:34} and Fact~\ref{fact:36}.
  Definition~\ref{def:47a} achieves the same in the same-visit case
  without changing stages.
\item \emph{Exclusion} (LMCS, proof of Thm.~9.11).  Forbid unfolding of
  the tile on which two conflicting tiles depend, and repeat.  On
  Example~\ref{ex:obstruction} this forbids unfolding $S_2$.  The
  paper's count of excluded tiles, $(\delta+p)N(n)$, does not cover this
  tile, which is an ancestor in the family rather than an x-special
  tile or one of its dependents.
\end{itemize}
\end{remark}

\subsection{Tree-shaped unfolding}\label{sec:treeunf}

Definition~\ref{def:unfold} builds the enlarged comb as a chain
$c_1c_2\cdots c_k$: each comb hangs below the leaf at which the previous
stage finished.  Example~\ref{ex:obstruction} fails because a second
visit to the same leaf continues in another direction (through the
constant~(4) to its second branch~(9)), and the chain has no room for
it.  The tree-shaped version copies the whole subtree instead, so every
direction keeps its place.

\begin{definition}[Tree-shaped unfolding]\label{def:treeunf}
\New{} Let $n$ be an $x$-node of $t$ and $m$ a non-lambda node strictly
below $n$.  Let $N$ and $M$ be the subterms of $t$ at $n$ and $m$, and
let $N^\circ$ be a copy of $N$ in which every lambda inside $N$ binds
fresh variables.  The \emph{tree-shaped unfolding of $t$ at $(n,m)$} is
\[
t'=t[m:=N^\circ].
\]
At the address of $m$ relative to $n$, the copy $N^\circ$ contains a
copy $M^\circ$ of $M$.  So $M$ survives, one level of $N$ deeper.
\end{definition}

So $t'$ is always a term: the copy has the shape of a subterm of $t$.
For $n=x_i\,A_1\dots A_k$ with $m$ in $A_j=\lambda\vec z.\,S_j$,
Definition~\ref{def:onestep} keeps the copies of $A_h$, $h\ne j$, as
here, but replaces the copy of $A_j$ by $\lambda\vec z'.\,z'_l\,\vec R'$.
The two definitions coincide iff $m$ lies directly below $\suc nj$.

\begin{definition}[Silent approach]\label{def:silent}
\New{} The pair $(n,m)$ is \emph{silent} if for every play $\pi$ of
$\Gm(t,E)$ and every position $\pi_a$ at $n$ the following holds.  If
$\pi_e$ is the first position after $\pi_a$ at $m$, then $\iv ae$
contains no $f$-move.
\end{definition}

\begin{theorem}[Tree-shaped unfolding is sound]\label{thm:treeunf}
\New{} Let $t$ solve $E$, let $n$ be a headless $x$-node
(Definition~\ref{def:head}), and let $(n,m)$ be silent.  Then the
tree-shaped unfolding $t'$ of $t$ at $(n,m)$ solves $E$.
\end{theorem}
\begin{proof}
We use the leftmost reduction sequences of the proof of
Theorem~\ref{thm:game}.  A \emph{state} is a closed term of ground
type.  A step $s\Rightarrow s'$ is either a $\beta$-step at the head, or
$f\,s_1\dots s_p\Rightarrow s_k$, where $\forall$ chooses $k$.  The
normal form of a state is determined by its tree of steps:
$\mathrm{nf}(s)=\mathrm{nf}(s')$ for a $\beta$-step,
$\mathrm{nf}(f\,\vec s)=f\,\mathrm{nf}(s_1)\dots\mathrm{nf}(s_p)$, and
$\mathrm{nf}(a)=a$.

Write $t=\lambda\vec x.\,T$, $t'=\lambda\vec x.\,T'$, and
$s^*=s[\vec w/\vec x]$, so that $t\,\vec w\to_\beta T^*$ and
$t'\,\vec w\to_\beta T'^*$.  Since $n$ is headless, every variable
occurring free in $N$ is one of the $x_i$: a $z$-variable bound above
$n$ would be a head carried by $n$, and the lambdas below constants are
dummies.  Hence $N^*$ and $N^{\circ*}$ are closed and $\alpha$-equivalent.

\emph{Labels.}  Label the occurrence of $M$ at $m$ in $T$, and follow it
through the steps: substitution copies labels.  A labelled occurrence
is an instance $M\theta$; it has ground type, so it is never in function
position.  By the correspondence in the proof of
Theorem~\ref{thm:game}, a state is labelled exactly when the
corresponding position is at $m$, and $f$-moves correspond to branch
steps.

\emph{Phase 1.}  Let $s\,R\,s'$ if $s'$ arises from $s$ by replacing
every labelled occurrence by $N^{\circ*}$.  Then $T^*\,R\,T'^*$.  If
$s\,R\,s'$ and $s$ is not itself labelled, then $s$ and $s'$ have the
same head and make the same step, with the same choice of $\forall$, and
$R$ holds again.  For a $\beta$-step this is because $N^{\circ*}$ is
closed, so replacing commutes with substitution:
$(b[a/y])'=b'[a'/y]$.

\emph{The switch.}  Fix a sequence of choices, and let $s_k$ be the
first labelled state, if there is one.  Then $s'_k=N^{\circ*}$.  Let
$\pi_e$ be the corresponding position at $m$.  By
Lemma~\ref{lem:ancestors} there is an earlier position $\pi_a$ at $n$,
with state $s_g$, $g<k$.  Its associated term is $N^*$, because $n$ is
headless.  $\pi_e$ is the first position at $m$ after $\pi_a$, since
$s_k$ is the first labelled state.  As $(n,m)$ is silent,
$s_g\Rightarrow\dots\Rightarrow s_k$ consists of $\beta$-steps only.
$N^{\circ*}$ equals $s_g$ up to $\alpha$ and labels, so it performs the
same $\beta$-steps: $N^{\circ*}\Rightarrow^*s''$, where $s''$ equals
$s_k$ up to $\alpha$ and labels.

\emph{Phase 2.}  From $s_k$ and $s''$ on, the two states are equal up
to $\alpha$ and labels, so they make the same steps forever.

So, for every sequence of choices, the steps from $T'^*$ are those from
$T^*$ with at most one block of $\beta$-steps inserted.  Hence
$\mathrm{nf}(T'^*)=\mathrm{nf}(T^*)=u$, and $t'$ solves $E$.
\end{proof}

\begin{remark}[About the hypotheses]
\begin{enumerate}
\item \emph{Silence is the real condition.}  The copy replays the part
  of the play from $n$ to $m$.  An $f$-move there would be performed
  twice.  In Example~\ref{ex:obstruction} the constant~(4) lies between
  $n$ and $m$, so $(n,m)$ is not silent, and indeed
  $t'\,w=f\,(f\,(f\,a\,d)\,d)\,d\ne u$.
\item \emph{Headlessness} is used only to make $N^*$ closed.  The
  argument should extend to every $x$-node by carrying the outer
  substitution along, but this is not written out.  If heads exist at
  all, a headless $x$-node whose successor binds a head always exists:
  take such an $x$-node of least depth.
\item \emph{Computer check} (not part of the proof).  All terms
  $\lambda x.\,T$ with $T$ of depth $\le3$ over $x$ unary, unary $z$,
  constants $f,g,a,d$, against $7$ arguments $w$ (linear and
  non-linear heads, with and without constants), and binary $x$ at
  depth~$2$ against $5$ arguments: in all $2766$ silent headless
  cases $t'\,\vec w$ and $t\,\vec w$ have the same normal form.  In all
  $1070$ non-silent cases they differ.  The $16$ silent cases with $n$
  not headless also agree.
\end{enumerate}
\end{remark}

\begin{remark}[Tree-shaped unfolding and compression]\label{rem:compress}
Tree-shaped unfolding never fails to give a term, and under silence it
is sound.  It does not compress, however.  The copy reproduces the
whole path from $\suc nj$ to $m$, including the $x$-nodes on it.  So
if $m$ was a head carried by such an $x$-node, the copy $M^\circ$ is
again a head of its copy, and the hydra measure does not drop.
Definition~\ref{def:onestep} is tree-shaped unfolding followed by a
\emph{cut} in the copy: everything strictly between $\lambda\vec z'$
and $M^\circ$ is removed, and variables of $M^\circ$ bound on the removed
part are bound by the originals again.  The open part of
\GapUnfolding{} is therefore the soundness of the cut.  It is a
$T_2$-like step on the copy, sound when no play needs the removed part.
Example~\ref{ex:obstruction} is the case where the second visit needs
it.
\end{remark}

%======================================================================
\section{The small solution theorem and decidability}\label{sec:thm37}
%======================================================================

\subsection{Counting nodes on a path}\label{sec:pathcount}

This subsection is purely combinatorial.  It replaces the pigeonhole
argument of [S, p.~53], see Section~\ref{sec:architecture}.

\begin{definition}[Strips]\label{def:strip}
Let $t\in T_A(\CE)$ and fix a path $P$ of $t$ from the root.  $L_x(P)$
and $L_z(P)$ are the numbers of $x$-nodes and of $z$-nodes on $P$.  A
\emph{strip} is a set of nodes on $P$ carrying the same variable with
the same binder.  $s_x(P)$ and $s_z(P)$ are the sizes of the largest
strips of $x$-nodes and of $z$-nodes on $P$.
\end{definition}

\begin{lemma}[Path counting]\label{lem:pathcount}
\New{} For every path $P$,
\[
L_x(P)\le q\,s_x(P),\qquad L_z(P)\le q\,s_z(P)\,L_x(P)\le q^2\,s_x(P)\,s_z(P),
\]
and for a longest path $P$,
\[
\depth(t)\le 2\bigl(L_x(P)+L_z(P)+\depth(u)+1\bigr).
\]
\end{lemma}
\begin{proof}
All $x$-nodes are bound by the root of $t$, which binds $k\le q$
variables, and each variable labels at most $s_x(P)$ nodes on $P$.  A
$z$-node on $P$ is bound by a successor of an $x$-node; that successor is
an ancestor of the $z$-node, so it lies on $P$ directly below an
$x$-node of $P$.  So there are at most $L_x(P)$ binders, each with at
most $q$ variables, each labelling at most $s_z(P)$ nodes on $P$.

For the depth: the nodes of $P$ are $x$- and $z$-nodes, constant nodes
of positive arity (at most $\depth(u)$, Lemma~\ref{lem:constants}), one
leaf, and one lambda node above each of these.
\end{proof}

\begin{proposition}[Strip bound]\label{prop:stripbound}
\begin{enumerate}
\item If $t$ is a terse solution, then $s_z(P)\le q\,\depth(u)+1$ for
  every path $P$ (Proposition~\ref{prop:stacking}).
\item \Open{} In the term produced by the construction of
  Section~\ref{sec:architecture} (exhaustive unfolding, then $T_1,T_2$),
  $s_x(P)\le K$ for some $K$ computable from $E$.  This is the content
  of the paper's ``if they are not end nodes and $k>q$ then some $N_1$
  can be unfolded'' (\GapCounting{}).
\end{enumerate}
\end{proposition}

\begin{corollary}[Depth from $s_x$]\label{cor:depthsz}
\New{} Let $\sigma=q\,\depth(u)+1$.  If $s_z(P)\le\sigma$ and $s_x(P)\le K$
for all paths, then
\[
L_x\le qK,\qquad L_z\le q^2\sigma K,\qquad
\depth(t)\le 2\bigl(qK+q^2\sigma K+\depth(u)+1\bigr).
\]
\end{corollary}
\begin{proof}
Lemma~\ref{lem:pathcount}.
\end{proof}

So everything reduces to one number: $s_x$, how many times the same
$x_i$ can stand one below another on a path.  This is the only place
where unfolding is needed, and by Proposition~\ref{prop:headless} only
for $x$-nodes of kinds (H) and (L) of Remark~\ref{rem:leftsx}.

\subsection{Heads and headless $x$-nodes}\label{sec:heads}

This subsection bounds the $x$-nodes on a path that neither carry a head
nor use their binder further down.  Only the remaining $x$-nodes need
unfolding.

\begin{definition}[Head]\label{def:head}
\New{} Let $n$ be an $x$-node of $t$.  A \emph{head carried by $n$} is a
$z$-node $m$ strictly below $n$ whose binder lies strictly above $n$.
The binder is then $\suc{n'}{h'}$ for an $x$-node $n'$ strictly above
$n$.  $n$ is \emph{headless} if it carries no head.
\end{definition}

\begin{lemma}[Confinement]\label{lem:confine}
\New{} Let $n$ be a headless $x$-node, $\pi\in\Gm(t,E)$ and $\pi_a$ a
position at $n$.  Then for every $c\ge a$:
\begin{enumerate}
\item $n_c$ lies in the subtree of $t$ at $n$ or in some $w_j$;
\item every entry of $\theta_c$ for a binder in the subtree at $n$ or in
  some $w_j$ is $((\suc{n_d}1,\dots),\theta_d)$ with $a\le d<c$.
\end{enumerate}
Consequently, $\pi$ has no other position at $n$, and every position of
$\pi$ at a successor of $n$ after $\pi_a$ is a child of $\pi_a$.
\end{lemma}
\begin{proof}
Induction on $c$.  For $c=a$: $n_a=n$, and the domain of $\theta_a$
consists of the binding nodes strictly above $n$
(Proposition~\ref{prop:tables}(1)).  So (2) is vacuous.  For the step,
look at the move from $\pi_c$ (Corollary~\ref{cor:jumps}).
\begin{itemize}
\item $\lambda$- and $f$-moves go to a successor, which stays in the
  subtree at $n$ or in $w_j$.  The table is unchanged.
\item $x$-jump at $n_c$: the target is the root of some $w_j$.  The new
  table has the single entry $((\suc{n_c}1,\dots),\theta_c)$, and
  $c\ge a$.
\item $z$-jump at a $z$-node $n_c$ in the subtree at $n$: its binder $b$
  is not strictly above $n$, since $n$ is headless.  As $n$ binds
  nothing, $b$ lies strictly below $n$.  By (2), $\theta_c(b)=
  ((\suc{n_d}1,\dots),\theta_d)$ with $d\ge a$, and $n_d$ is a $y$-node.
  The target $\suc{n_d}g$ lies in $w_j$.  The new table is $\theta_d$,
  which satisfies (2) by induction, possibly extended by an entry with
  inner table $\theta_c$.
\item $y$-jump at $n_c$ in $w_j$, bound by the root of $w_j$: by (2) the
  entry is $((\suc{n_d}1,\dots),\theta_d)$ with $d\ge a$.  Here $n_d$ is
  an $x$-node, in the subtree at $n$ by (1).  The target $\suc{n_d}h$ is
  in that subtree.  The new table is $\theta_d$ extended by an entry
  with inner table $\theta_c$.
\item $v$-jump at $n_c$ in $w_j$: likewise, the target is a dummy lambda
  below a $z$-node $n_d$ with $d\ge a$, so it is in the subtree at $n$,
  and the new table is $\theta_d$.
\end{itemize}
For the consequences: a later position at $n$ would follow a position at
the lambda node directly above $n$ (Lemma~\ref{lem:pred}).  That node
is outside the subtree at $n$, which contradicts (1).  The parent of a
position at $\suc nh$ is a position at $n$ (Definition~\ref{def:parent}),
so it is $\pi_a$.
\end{proof}

\begin{lemma}[Witness of non-redundancy]\label{lem:witness}
\New{} Let $t$ solve $E$, and let $\gamma=(n,\suc nh)$ be an end path
($x$-path with $p=1$) that is not redundant in $\Gm(t,E)$.  Then some
play $\pi$ has a position $\pi_a$ at $n$ such that either
\begin{enumerate}[(W1)]
\item $\pi_a$ has no child at $\suc nh$; or
\item $\iv ab$ contributes, where $\pi_b$ is the first child of
  $\pi_a$ at $\suc nh$.
\end{enumerate}
\end{lemma}
\begin{proof}
Suppose that no play has such a position.  We show by induction on
$l-i$ that $\gamma$ is redundant in $\iv il$ for every play.  If no
$j_1\in[i,l)$ has $n_{j_1}=n$, this holds.  Otherwise let $j_1$ be the
least such index and $\pi_j$ the first child of $\pi_{j_1}$ at $\suc nh$.
It exists since (W1) fails.  Then $\pi_{j_1},\pi_j$ is a chain for
$\pi_j$, and $\iv{j_1}j$ does not contribute since (W2) fails.  By
induction, $\gamma$ is redundant in $\iv jl$, because $j>i$.  So
Definition~\ref{def:redundant} holds with this $j$.
\end{proof}

\begin{proposition}[Headless end $x$-nodes]\label{prop:headless}
\New{} Let $t$ be a terse solution and $P$ a path of $t$ from the root.
For an $x$-node $n$ on $P$, let $\suc nh$ be its successor on $P$.  Call
$n$ \emph{free} if $n$ is headless and no variable node strictly below
$\suc nh$ is bound by $\suc nh$.  Then $P$ has at most $|u|$ free
$x$-nodes.
\end{proposition}
\begin{proof}
Let $n_1,\dots,n_K$ be the free $x$-nodes of $P$, from top to bottom,
with successors $\suc{n_k}{h_k}$ on $P$.  By Lemma~\ref{lem:endpath5},
$\gamma_k=(n_k,\suc{n_k}{h_k})$ is an end path.  Since $t$ is terse,
$\gamma_k$ is not redundant.  Lemma~\ref{lem:witness} gives a play
$\pi^k$ and a position $\pi^k_{a_k}$ at $n_k$ with (W1) or (W2).

\emph{Claim.}  Let $k<k'$ and let $\pi$ be a play with a position
$\pi_{a'}$ at $n_{k'}$.  Then $\pi$ has a position $\pi_a$ at $n_k$ with
$a<a'$, and $\pi_a$ has a child $\pi_e$ at $\suc{n_k}{h_k}$ with
$e<a'$.  Indeed, $\pi_{a'-1}$ is at the lambda node directly above
$n_{k'}$ (Lemma~\ref{lem:pred}), which is $\suc{n_k}{h_k}$ or lies below
it.  The chain for $\pi_{a'-1}$ from $n_k$ (Lemma~\ref{lem:ancestors})
starts with a position at $n_k$ followed by one of its children at
$\suc{n_k}{h_k}$, both at indices at most $a'-1$.  By Lemma~\ref{lem:confine}, $\pi_a$ is the only position of $\pi$
at $n_k$.

\emph{(W1)-indices.}  If $k$ satisfies (W1) in $\pi^k$, then by the
claim $\pi^k$ visits no $n_{k'}$ with $k'>k$.  So $k$ is the largest
index visited in $\pi^k$, and distinct (W1)-indices have distinct plays.
There are at most $\br(u)$ of them (Fact~\ref{fact:branches}).

\emph{(W2)-indices.}  If $k$ satisfies (W2) in $\pi^k$, then
$\iv{a_k}{b_k}$ contains an $f$-move at some $c_k\in[a_k,b_k)$
(Lemma~\ref{lem:states}).  We show that $k\mapsto c_k$ is injective as a
map into the tree of plays.  Suppose $c_k=c_{k'}=c$ with $k<k'$, so
$\pi^k$ and $\pi^{k'}$ agree up to $c$.  Then $a_{k'}\le c$, so
$\pi^k$ has the position $\pi^k_{a_{k'}}$ at $n_{k'}$.  By the claim
and the uniqueness of $\pi^k_{a_k}$, the first child $b_k$ of
$\pi^k_{a_k}$ at $\suc{n_k}{h_k}$ comes before $a_{k'}\le c$.  This
contradicts $c<b_k$.  So there are at most $|u|-\br(u)$ (W2)-indices
(Lemma~\ref{lem:global}).

Every $k$ is a (W1)- or a (W2)-index, so $K\le|u|$.
\end{proof}

\begin{remark}[What is left for $s_x$]\label{rem:leftsx}
The proposition does not need the $n_k$ to carry the same variable.
On a path $P$ of a terse solution, every $x$-node except at most $|u|$
is therefore of one of two kinds:
\begin{itemize}
\item[(H)] it carries a head; or
\item[(L)] it is headless, but its successor $\suc nh$ on $P$ binds a
  $z$-node below $\suc nh$.  That $z$-node is not below any $x$-node of
  $P$ further down, since that $x$-node would then carry a head.  So it
  lies on $P$ before the next $x$-node, or on a side branch.
\end{itemize}
Heads are what makes play leave and re-enter a subtree; without them
the intervals $\iv{a_k}{b_k}$ are disjoint, which is what the proof
counts.  Kinds (H) and (L) are exactly the configurations that
unfolding has to remove.  For (H) this is the hydra picture: an
unfolding at a head removes it, at the price of new heads higher up.
\end{remark}

\begin{lemma}[$T_1$ and $T_2$ create no heads]\label{lem:monoheads}
\New{} Let $t\,T_1\,t'$ or $t\,T_2\,t'$.  Every $x$-node of $t'$ comes
from an $x$-node of $t$, and if it is free on a path of $t'$, it was
free on the corresponding path of $t$ already, and conversely: an
$x$-node of $t'$ of kind (H) or (L) comes from one of kind (H) or (L)
in $t$.
\end{lemma}
\begin{proof}
$T_1$ replaces the subtree at a node by a leaf $b$.  $T_2$ deletes the
nodes of $U$ and moves $S$ up along $\varrho$
(Lemma~\ref{lem:T2sim}).  In both cases the surviving nodes keep their
labels, successor indices and binders, and one surviving node is below
another in $t'$ iff it is below it in $t$.  A head of $t'$ (a $z$-node
$m$ below an $x$-node $n$, with binder above $n$) is therefore a head
of $t$, and a $z$-node below $\suc nh$ bound by $\suc nh$ in $t'$ is one
in $t$.  Paths of $t'$ map to paths of $t$ with the same successor
indices at surviving nodes.
\end{proof}

\subsection{The theorem}

\begin{theorem}[S, Thm.~37 at order $\le5$, weak form]\label{thm:small}
\Open{} Let $E$ have order at most $5$, and let $q$ be as in
Definition~\ref{def:kinds}.  There is a number $D_E$, computable from
$E$, such that if $E$ has a (canonical) solution, it has one of depth at
most $D_E$.
\end{theorem}

\begin{remark}
\Fixed{} Only some computable bound is needed for decidability
(Corollary~\ref{cor:decidable}).  The argument below gives
$D_E=2\bigl(qK+q^2\sigma K+\depth(u)+1\bigr)$ from
Corollary~\ref{cor:depthsz}, with $K$ the bound on $x$-strips.  The
paper's sharper bound ($4p^2q^{2p-2}|u|$, i.e.\ $16q^2|u|$ for $p=2$) is
not pursued.
\end{remark}

\subsection{Proof architecture}\label{sec:architecture}

The proof goes from a terse solution $t$ to a terse solution $t^*$ whose
depth is bounded by a number computable from $E$.  It has five steps.
Each step lists the results it uses and the gaps it depends on.

\paragraph{Step 1: properties of a terse solution.}
Let $t$ be a terse solution (Definition~\ref{def:terse}); one exists if
$E$ has a solution (Lemma~\ref{lem:terse}).  Then:
\begin{enumerate}[(a)]
\item $t$ consists of $x$-nodes (level~$1$), $z$-nodes (level~$2$),
  constant nodes and lambda nodes, with branching and binder length at
  most $q$ (Corollary~\ref{cor:binders});
\item $\Gm(t,E)$ has $\br(u)\le|u|$ plays (Fact~\ref{fact:branches});
\item every node except the $b$-leaves and the dummy lambdas above them
  is visited (Corollary~\ref{cor:visited});
\item every path has at most $\depth(u)$ constants of positive arity
  (Lemma~\ref{lem:constants});
\item every end path is non-redundant, since $T_2$ does not apply; in
  particular $z$-strips have at most $q\,\depth(u)+1$ nodes
  (Proposition~\ref{prop:stacking}).
\end{enumerate}
\emph{Gaps:} none; existence is Lemma~\ref{lem:terse}.

\paragraph{Step 2: the hyper-partition tree.}
Partition every play into stages (Definition~\ref{def:stages}); every
stage is a constant stage, an $x$-stage or a $z$-stage
(Lemma~\ref{lem:stages5}).  Glue the plays into the hyper-partition tree
$\mathcal H$ (Definition~\ref{def:hyper}).  It is a tree, its partition
edges are well defined, and every separator contains a contributing
node (Proposition~\ref{prop:hyper}).  Gluing the combs along the edges
should give back the visited part of $t$.
\emph{Gaps:} that last claim is \GapPartition{} (partition tree).

\paragraph{Step 3: unfolding.}
Unfold according to the regime of Definition~\ref{def:regime}: always at
the $\prec$-least unfoldable $z$-node, i.e.\ from the bottom right of the
term to the left, by one-step unfoldings (Definition~\ref{def:onestep}),
with non-contribution read over $\mathcal H$ (Definition~\ref{def:47a}).
An unfolding takes a non-contributing $x$-stage at a node $n$ and at most
$q$ non-contributing $z$-stages bound by one successor of $n$; the
enlarged comb has depth at most $2(q+1)$
(Lemma~\ref{lem:cembshape}).  Separators are never unfolded
(Proposition~\ref{prop:hyper}(4)).  Unfolding must preserve solutions:
this is proved for linear heads (Theorem~\ref{thm:linunfold}), and for
non-linear heads the obstruction of Example~\ref{ex:obstruction} is
excluded in the same-visit case (Proposition~\ref{prop:reentry}).
\emph{Gaps:} soundness in general, the game of the unfolded term
(Fact~\ref{fact:36}) and the cross-visit case are \GapUnfolding{}.

\paragraph{Step 4: shrinking.}
Apply $T_1$ and $T_2$ exhaustively to the unfolded term; call the result
$t^*$.  It is terse by construction, so Step~1 applies to it.  By Lemma~\ref{lem:endpath5}, $T_2$ either replaces a $z$-node by
the body of one of its arguments, or replaces an $x$-node $n$ by the body
below a sequence of directly nested $z$-nodes bound at a successor of
$n$.  Both transformations preserve solutions and make the term smaller
(Propositions~\ref{prop:T1} and~\ref{prop:T2}).
\emph{Gaps:} termination of the whole
construction, and whether Step~3 must be repeated after Step~4, are
\GapCounting{}(i).

\paragraph{Step 5: the depth bound.}
By Lemma~\ref{lem:pathcount} the depth of $t^*$ is bounded once its
strips are bounded.
\begin{enumerate}[(a)]
\item $z$-strips: $s_z\le q\,\depth(u)+1$, since $z$-nodes are end
  (Proposition~\ref{prop:stripbound}(1)), which applies because $t^*$
  is terse.
\item $x$-strips: $s_x\le K$ for a computable $K$
  (Proposition~\ref{prop:stripbound}(2)).  On any path at most $|u|$
  $x$-nodes are free, i.e.\ headless and not binding below their
  successor (Proposition~\ref{prop:headless}, proved).  The others carry
  a head (H) or use their binder locally (L)
  (Remark~\ref{rem:leftsx}).  These have been removed by the
  exhaustive unfolding of Step~3, and Step~4 does not create new ones
  (Lemma~\ref{lem:monoheads}).  So in $t^*$ every $x$-node is free,
  and $L_x\le|u|$ on every path; in particular $s_x\le|u|$.
\end{enumerate}
Then Corollary~\ref{cor:depthsz} bounds the depth of $t^*$, and
Corollary~\ref{cor:decidable} gives decidability.
\emph{Gaps:} (a) is proved.  In (b), the bound on free $x$-nodes and
the stability under Step~4 are proved; that the unfolding of Step~3
leaves no $x$-node of kind (H) or (L) is \GapCounting{}(ii), the main
open point of the proof.

\renewcommand{\gapname}{[Gap-Counting]}
\begin{gapenv}\phantomsection\label{gap:count}
\Open{}
\begin{enumerate}[(i)]
\item Termination of the construction: unfolding enlarges the term and
  $T_1,T_2$ shrink it, and they may have to be interleaved.
\item The bound $s_x\le K$ (Proposition~\ref{prop:stripbound}(2)).
  Free $x$-nodes are at most $|u|$ per path
  (Proposition~\ref{prop:headless}), and $T_1,T_2$ create no new
  $x$-nodes of kinds (H) and (L) (Lemma~\ref{lem:monoheads}).  Open:
  every $x$-node of kind (H) or (L) (Remark~\ref{rem:leftsx}) admits an
  unfolding that removes it, or is charged to a contribution as in
  Proposition~\ref{prop:headless}.
\end{enumerate}
\end{gapenv}

\begin{corollary}[Decidability]\label{cor:decidable}
\Paper{} \Coq{} Assume Theorem~\ref{thm:small}, or just a computable
bound on $s_x$ (Proposition~\ref{prop:stripbound}(2)).  Then canonical
solvability of interpolation equations of order at most~$5$ under (R1),
(R2) is decidable, and so is every matching problem that reduces to
such an equation.
\end{corollary}
\begin{proof}
A matching equation reduces to the canonical solvability of an
interpolation equation [S, Prop.~5], with constants in $\CE$
(Lemma~\ref{lem:reduces}).  Up to $\alpha$, there are finitely many lnf
terms of bounded depth over $\CE$, since names can be drawn from a
fixed finite set.  Enumerate them and check each by normalisation.  In
Coq this is \texttt{decidable\_interpolation} (\texttt{Small.v}), with
\texttt{Thm37} as a hypothesis.
\end{proof}

%======================================================================
\section{Removing the restrictions [S, \S6.5]}\label{sec:restrictions}
%======================================================================

Summary only; not re-checked.
\begin{itemize}
\item \emph{Several base types.} Nodes carry a result type, which can
  change only at constants, so at most $\depth(u)$ times per play.
  $T_1$ uses one constant $b_B$ per base type, and $T_2$ and unfolding
  are restricted to paths of constant result type.  The bound gets
  $|u|^2$ in place of $|u|$.
\item \emph{Goals with bound variables.} Tables get entries for binders
  below constants, mapping their variables to fresh constants.  Theorem~13
  must be re-proved for the amended game.
\end{itemize}

%======================================================================
\section{Summary of status}\label{sec:summary}
%======================================================================

\begin{center}
\small
\begin{tabular}{lll}
Here & Paper & Status\\\hline
Prop.~\ref{prop:tables}, Thm.~\ref{thm:game}, Prop.~\ref{prop:distinct} & Prop.~12, Thm.~13, Prop.~14 & proved (paper); Coq axioms\\
Props.~\ref{prop:parent}--\ref{prop:agree} & Prop.~15--19 & proved (paper); Coq axioms\\
Lemmas~\ref{lem:pred}--\ref{lem:alt} & implicit & proved; partly Coq\\
Prop.~\ref{prop:between}, \ref{prop:24} & Props.~23--24 & proved (paper)\\
Prop.~\ref{prop:24star}, Cor.~\ref{cor:24prime} & --- & new, proved\\
Facts~\ref{fact:25}, \ref{fact:26} & Facts~25--26 & proved\\
Fact~\ref{fact:27} & Fact~27 & sketch\\
Prop.~\ref{prop:29} & Prop.~29 & corrected (greedy chain)\\
Prop.~\ref{prop:T1} & Fact~30 & proved; Coq\\
Lemma~\ref{lem:T2sim}, Prop.~\ref{prop:T2} & Prop.~31 & new simulation lemma, proved\\
Prop.~\ref{prop:stacking} & Prop.~32 & corrected bound, proved\\
Fact~\ref{fact:33}, Prop.~\ref{prop:34} & Fact~33, Prop.~34 & proved (Coq: uniqueness)\\
Def.~\ref{def:ptree2} & Def.~46 & corrected reading, \GapPartition{}\\
Def.~\ref{def:hyper}, Prop.~\ref{prop:hyper} & p.~53 & new, proved\\
Lemma~\ref{lem:shape}, Cor.~\ref{cor:binders}, \ref{cor:jumps} & --- & new, proved (order-5 shape)\\
Lemma~\ref{lem:endpath5} & Def.~39 & new, proved\\
Lemma~\ref{lem:stages5} & Prop.~34 & new, proved\\
Lemma~\ref{lem:cembshape} & p.~53 & new, proved\\
Thm.~\ref{thm:linunfold} & --- & new; Coq (raw terms)\\
Ex.~\ref{ex:obstruction} & Def.~47 & counterexample\\
Prop.~\ref{prop:reentry}, Cor.~\ref{cor:crossvisit} & --- & new; partial, \GapUnfolding{}\\
Fact~\ref{fact:36} & Fact~36 & open\\
Def.~\ref{def:treeunf}, Thm.~\ref{thm:treeunf} & --- & new; sound under silence, proved\\
Lemma~\ref{lem:pathcount}, Cor.~\ref{cor:depthsz} & p.~53 & new, proved (replaces pigeonhole)\\
Prop.~\ref{prop:stripbound} & p.~53 & (1) proved, (2) open\\
Lemmas~\ref{lem:confine}, \ref{lem:witness}, \ref{lem:monoheads}, Prop.~\ref{prop:headless} & --- & new, proved\\
Thm.~\ref{thm:small} & Thm.~37 & weak form (computable bound); open\\
Cor.~\ref{cor:decidable} & \S6 & proved from Thm.~37; Coq
\end{tabular}
\end{center}

\paragraph{Corrections adopted.}
\begin{enumerate}
\item End paths: ``below $m_{2p}$'' read strictly
  (Definition~\ref{def:endpath}).
\item Contribution of intervals not ending at a lambda node
  (Definition~\ref{def:contrib}).
\item Prop.~29 and stages use the greedy recurrent chain
  (Proposition~\ref{prop:29}).
\item Stage cut: the first $p$ with $\pi_{j_{2p}+2}$ a child of $\pi_i$
  (Definition~\ref{def:stages}).
\item Partition-tree edges via parents (Definition~\ref{def:ptree2}).
\item Prop.~32 gives $p\le q\,\depth(u)+1$, and $p\le\depth(u)+1$ when the
  strip uses one successor index throughout.
\item Def.~47(a) read over all plays (Definition~\ref{def:47a}).  As
  written, Def.~47 can produce a non-term
  (Example~\ref{ex:obstruction}).
\item Thm.~37: $q$ also bounds the arities of the constants of $u$.
\item $T_1$ also replaces unvisited 0-ary constants other than $b$, with
  size measure $(|t|,\nu,\kappa)$, so that in a terse solution exactly
  the $b$-leaves (with their lambdas) are unvisited
  (Corollary~\ref{cor:visited}).
\item The least-$j$ pigeonhole of Thm.~37 is replaced by
  Lemma~\ref{lem:pathcount} ($L_x\le q\,s_x$, $L_z\le q\,s_z\,L_x$) and a
  bound on strips (Proposition~\ref{prop:stripbound}).
\end{enumerate}

\paragraph{Where the gaps are.}
\begin{center}
\small
\begin{tabular}{llll}
Gap & Topic & Introduced in & Page\\\hline
\GapPartition & Partition tree & \S\ref{sec:partition}, after Def.~\ref{def:ptree2} & \pageref{gap:ptree}\\
\GapUnfolding & Unfolding & \S\ref{sec:nonlinear}, after Cor.~\ref{cor:crossvisit} & \pageref{gap:unfold}\\
\GapCounting & Counting & \S\ref{sec:architecture}, after Step~5 & \pageref{gap:count}
\end{tabular}
\end{center}

\paragraph{Open, in order of importance.}
\GapUnfolding{} (\S\ref{sec:nonlinear}; unfolding: game, soundness,
the remaining cross-visit case), \GapCounting{}
(\S\ref{sec:architecture}; counting), \GapPartition{}
(\S\ref{sec:partition}; partition tree).  [Gap-Stacking] is closed by
Proposition~\ref{prop:stacking}, and [Gap-$T_2$] by
Lemma~\ref{lem:T2sim}.

\end{document}
